% !TEX root = 第九章_函数项级数、幂级数、Fourier级数(答案版).tex
\setcounter{chapter}{8}
\pagestyle{fancy}

\chapter{函数项级数,幂级数,Fourier 级数}


\section{一致收敛性}

判断函数列或函数项级数一致收敛性的方法如下.

判断一致收敛性:
\begin{enumerate}
    \item 定义(三分法);
    \item 验证 $\displaystyle\lim\limits_{n\to\infty}M_n=0$,其中
          $M_n=\displaystyle\sup\limits_{x\in I}|f_n(x)-f(x)|$.验证时通常计算极值点,或直接放缩为趋于 $0$ 的常数列;
    \item Cauchy 收敛准则;
    \item Weierstrass,Abel,Dirichlet 判别法;
    \item Dini 定理;
    \item 利用等度连续.
\end{enumerate}

判别非一致收敛性:
\begin{enumerate}
    \item 通项不一致趋于 $0$;
    \item 端点法(一般用于有限区间):若 $f_n\in C\lbrack a,b\rbrack$,且 $\displaystyle\sum\limits_{n=1}^{\infty}f_n(x)$ 在 $(a,b)$ 上一致收敛,则必可延拓为在 $\lbrack a,b\rbrack$ 上一致收敛.因此若级数在端点发散,则在 $(a,b)$ 上非一致收敛;
    \item 利用一致收敛传递连续性:若连续函数列的极限函数在 $I$ 上不连续,则在 $I$ 上非一致收敛;
    \item Cauchy 收敛准则,其优点在于不必考虑或计算极限函数.
\end{enumerate}

\begin{remark}
    函数列 $\{f_n(x)\}$ 的问题可通过考虑 $\displaystyle\sum\limits_n[f_{n+1}(x)-f_n(x)]$ 转化为函数项级数问题,而函数项级数可通过考虑前 $n$ 项和 $\{s_n(x)\}$ 转化为函数列问题.
\end{remark}

%例题9.1.1
\begin{example}
    试判断下列函数列在指定区间上的一致收敛性:
    \begin{enumerate}
        \item $f_n(x)=x^n-x^{n+1}$, $x\in\lbrack 0,1\rbrack$;
        \item $f_n(x)=\dfrac{x}{1+nx^2}$, $x\in(-\infty,+\infty)$;
        \item $f_n(x)=n^\alpha e^{-nx}$, $x\in\lbrack 0,1\rbrack$.
    \end{enumerate}
\end{example}
\exampleblank

\begin{solution}
    \leavevmode
    \begin{enumerate}
        \item 极限函数为零,且
              \[
                  \max_{0\leq x\leq1}x^n(1-x)
                  =\dfrac{1}{n+1}\left(\dfrac{n}{n+1}\right)^n\longrightarrow0,
              \]
              故一致收敛.
        \item 极限函数为零,且
              \[
                  \sup_{x\in\mathbb{R}}\dfrac{|x|}{1+nx^2}=\dfrac{1}{2\sqrt n}\longrightarrow0,
              \]
              故一致收敛.
        \item 当 $\alpha<0$ 时,$0\leq f_n(x)\leq n^\alpha\to0$,故一致收敛于零.当 $\alpha=0$ 时,逐点极限在 $x=0$ 处为 $1$,在 $x>0$ 处为零,不连续,故不一致收敛;当 $\alpha>0$ 时,$f_n(0)=n^\alpha$,连逐点有限收敛也不成立.
    \end{enumerate}
\end{solution}

%例题9.1.2
\begin{example}
    试判断下列函数列在指定区间上的一致收敛性:
    \begin{enumerate}
        \item $f_n(x)=\sin\dfrac{1+nx}{2n}$, $x\in(-\infty,+\infty)$;
        \item $f_n(x)=n\sin\dfrac{1}{nx}$, $x\in\lbrack 1,+\infty)$;
        \item $f_n(x)=\dfrac{\ln(nx)}{nx^2}$, $x\in(1,+\infty)$;
        \item $f_n(x)=\dfrac{n}{x}\ln\left(1+\dfrac{x}{n}\right)$, $x\in(0,10)$.
    \end{enumerate}
\end{example}
\exampleblank

\begin{solution}
    \leavevmode
    \begin{enumerate}
        \item 极限为 $\sin(x/2)$,且由正弦函数的 Lipschitz 性,
              \[
                  \sup_{x\in\mathbb{R}}\left|f_n(x)-\sin\dfrac{x}{2}\right|\leq\dfrac{1}{2n},
              \]
              故一致收敛.
        \item 极限为 $1/x$.由 $|\sin t-t|\leq|t|^3/6$,
              \[
                  \sup_{x\geq1}\left|n\sin\dfrac{1}{nx}-\dfrac{1}{x}\right|
                  \leq\dfrac{1}{6n^2},
              \]
              故一致收敛.
        \item 极限为零.充分大时 $x\mapsto\ln(nx)/x^2$ 在 $x>1$ 上递减,故
              \[
                  \sup_{x>1}|f_n(x)|=\dfrac{\ln n}{n}\longrightarrow0.
              \]
        \item 极限为 $1$,且 $0<\ln(1+t)-t+t^2/2<t^2/2$ 给出
              \[
                  \sup_{0<x<10}|f_n(x)-1|\leq\dfrac{5}{n}\longrightarrow0.
              \]
    \end{enumerate}
\end{solution}

%例题9.1.3
\begin{example}
    证明下列函数列在指定区间上的非一致收敛性:
    \begin{enumerate}
        \item $f_n(x)=x^n-x^{2n}$, $x\in\lbrack 0,1\rbrack$;
        \item $f_n(x)=x^n+x^{2n}-2x^{3n}$, $x\in\lbrack 0,1\rbrack$;
        \item $f_n(x)=nx(1-x)^n$, $x\in\lbrack 0,1\rbrack$;
        \item $f_n(x)=n\left(\sqrt{x+\dfrac{1}{n}}-\sqrt{x}\right)$, $x\in(0,+\infty)$.
    \end{enumerate}
\end{example}
\exampleblank

\begin{solution}
    四列均逐点趋于零,但可分别选取
    \[
        x_n^{(1)}=2^{-1/n},\qquad x_n^{(2)}=2^{-1/n},\qquad
        x_n^{(3)}=\dfrac{1}{n+1},\qquad x_n^{(4)}=\dfrac{1}{n^2}.
    \]
    此时前两列的函数值分别恒为 $1/4$ 与 $1/2$;第三列趋于 $e^{-1}$;第四列为
    \[
        \dfrac{1}{\sqrt{n^{-2}+n^{-1}}+n^{-1}},
    \]
    甚至趋于 $+\infty$.因此都不一致收敛.
\end{solution}

%例题9.1.4
\begin{example}
    试判断下列级数 $\displaystyle I=\sum\limits_{n=1}^{\infty}f_n(x)$ 在指定区间上的一致收敛性:
    \begin{enumerate}
        \item $\displaystyle\sum\limits_{n=1}^{\infty}x^\alpha e^{-nx^2}$, $x\in(0,+\infty)$;
        \item $\displaystyle\sum\limits_{n=1}^{\infty}e^{-(x-n)^2}$, $x\in\lbrack a,b\rbrack$;
        \item $\displaystyle\sum\limits_{n=1}^{\infty}\dfrac{x^n}{4^n}\ln\left(1+\dfrac{x^2}{n^2}\right)$, $x\in\lbrack -4,4\rbrack$;
        \item $\displaystyle\sum\limits_{n=1}^{\infty}\dfrac{x}{n^p+n^q x^2}$, $x\in\lbrack 0,+\infty)$.
    \end{enumerate}
\end{example}
\exampleblank

\begin{solution}
    \leavevmode
    \begin{enumerate}
        \item 当 $\alpha>0$ 时,
              \[
                  \sup_{x>0}x^\alpha e^{-nx^2}=C_\alpha n^{-\alpha/2}.
              \]
              因此 $\alpha>2$ 时由 Weierstrass 判别一致收敛.若 $\alpha\leq2$,在尾项中取 $x=N^{-1/2}$ 并考察 $N\leq n\leq2N$,尾和不趋于零,故不一致收敛.
        \item 设 $C=\max\{|a|,|b|\}$.充分大时
              $e^{-(x-n)^2}\leq e^{-(n-C)^2}$,故一致收敛.
        \item 对 $|x|\leq4$,
              \[
                  \left|\dfrac{x^n}{4^n}\ln\left(1+\dfrac{x^2}{n^2}\right)\right|
                  \leq\dfrac{16}{n^2},
              \]
              故一致绝对收敛.
        \item
              \[
                  \sup_{x\geq0}\dfrac{x}{n^p+n^qx^2}=\dfrac{1}{2n^{(p+q)/2}}.
              \]
              因而 $p+q>2$ 时一致收敛;反之取 $x=N^{(p-q)/2}$ 并对
              $N\leq n\leq2N$ 求和,尾和不趋于零.故恰在 $p+q>2$ 时一致收敛.
    \end{enumerate}
\end{solution}

%例题9.1.5
\begin{example}
    试判断下列级数在指定区间上的一致收敛性:
    \begin{enumerate}
        \item $\displaystyle\sum\limits_{n=1}^{\infty}\dfrac{2nx}{1+n^\alpha x}$, $-\infty<x<+\infty$, $\alpha>\varphi$;
        \item $\displaystyle\sum\limits_{n=1}^{\infty}\dfrac{x\sin nx}{\sqrt{1+n^2(1+nx^2)}}$, $x\in(-\infty,+\infty)$;
        \item $\displaystyle\sum\limits_{n=1}^{\infty}\dfrac{x^2}{1+n^2x}\sin\dfrac{n}{x}$, $x\in(0,1)$;
        \item $\displaystyle\sum\limits_{n=1}^{\infty}\dfrac{(-1)^{n+1}}{2n+\sin x}$, $x\in(-\infty,+\infty)$.
    \end{enumerate}
\end{example}
\exampleblank

\begin{solution}
    \leavevmode
    \begin{enumerate}
        \item 当前题面在 $x=-n^{-\alpha}$ 处分母为零,故函数甚至未在整个实轴上定义,无法讨论该区间上的一致收敛性;区间或分母符号需要核对.
        \item
              \[
                  \left|\dfrac{x\sin nx}{\sqrt{1+n^2(1+nx^2)}}\right|
                  \leq\dfrac{|x|}{n\sqrt{1+nx^2}}\leq\dfrac{1}{n^{3/2}},
              \]
              故一致绝对收敛.
        \item 对 $0<x<1$,
              \[
                  \left|\dfrac{x^2}{1+n^2x}\sin\dfrac{n}{x}\right|\leq\dfrac{x}{n^2}\leq\dfrac{1}{n^2},
              \]
              故一致绝对收敛.
        \item $b_n(x)=1/(2n+\sin x)$ 关于 $n$ 递减,且
              $\sup_xb_n(x)\leq1/(2n-1)\to0$,故由一致 Leibniz 判别法一致收敛.
    \end{enumerate}
\end{solution}

%例题9.1.6
\begin{example}
    试判断下列级数在指定区间上的非一致收敛性:
    \begin{enumerate}
        \item $\displaystyle\sum\limits_{n=1}^{\infty}\dfrac{x^3}{(1+x^3)^n}$, $x\in(0,+\infty)$;
        \item $\displaystyle\sum\limits_{n=1}^{\infty}\dfrac{x}{1+n^3x^3}$, $x\in(0,+\infty)$;
        \item $\displaystyle\sum\limits_{n=1}^{\infty}nx^2e^{-nx}$, $x\in(0,+\infty)$;
        \item $\displaystyle\sum\limits_{n=1}^{\infty}\dfrac{x}{x^2-nx+n^2}$, $x\in(1,+\infty)$.
    \end{enumerate}
\end{example}
\exampleblank

\begin{solution}
    \leavevmode
    \begin{enumerate}
        \item 几何级数的余项为 $(1+x^3)^{-N}$,其在 $x>0$ 上的上确界为 $1$,故不一致收敛.
        \item 取 $x=1/N$,则从 $n=N$ 到 $2N$ 的尾和为
              \[
                  \displaystyle\sum\limits_{n=N}^{2N}\dfrac{N^{-1}}{1+(n/N)^3}\geq c>0.
              \]
        \item 仍取 $x=1/N$,有
              \[
                  \displaystyle\sum\limits_{n=N}^{2N}nx^2e^{-nx}geq c>0.
              \]
        \item 取 $x=N$,对 $N\leq n\leq2N$,每项均为 $N^{-1}$ 量级,故相应尾和有正下界.
    \end{enumerate}
    四者均违反一致 Cauchy 准则.
\end{solution}

%例题9.1.7
\begin{example}
    试用 Cauchy 收敛准则判别下列级数在指定区间上的非一致收敛性:
    \begin{enumerate}
        \item $\displaystyle\sum\limits_{n=1}^{\infty}\dfrac{n}{x}e^{-n^2/x}$, $x\in(0,+\infty)$;
        \item $\displaystyle\sum\limits_{n=1}^{\infty}\dfrac{\sin nx}{n}$, $x\in\lbrack 0,2\pi\rbrack$;
        \item $\displaystyle\sum\limits_{n=1}^{\infty}\dfrac{\sin nx}{e^{n^2x}}$, $x\in(0,+\infty)$;
        \item $\displaystyle\sum\limits_{n=1}^{\infty}\dfrac{\sqrt{\pi}}{n^2}\sin\dfrac{\pi}{n^2}$, $x\in(0,+\infty)$.
    \end{enumerate}
\end{example}
\exampleblank

\begin{solution}
    \leavevmode
    \begin{enumerate}
        \item 取 $x=N^2$ 并令 $N\leq n\leq2N$,则尾和
              $\displaystyle\sum(n/N^2)e^{-n^2/N^2}$ 有固定正下界.
        \item 取 $x=1/N$,则 $N\leq n\leq2N$ 时 $nx\in[1,2]$,相应尾和有固定正下界.
        \item 取 $x=1/N^2$,利用 $\sin(nx)\sim nx$,从 $N$ 到 $2N$ 的尾和仍有固定正下界.
        \item 当前级数与 $x$ 无关,并且
              \[
                  \left|\dfrac{\sqrt\pi}{n^2}\sin\dfrac{\pi}{n^2}\right|
                  \leq\dfrac{\pi^{3/2}}{n^4}.
              \]
              因而它实际上一致绝对收敛,与题目要求``非一致收敛''矛盾,题面应有漏项.
    \end{enumerate}
\end{solution}

%例题9.1.8
\begin{example}
    试用端点法判别下列级数在指定区间上的非一致收敛性:
    \begin{enumerate}
        \item $\displaystyle\sum\limits_{n=1}^{\infty}e^{-nx}\cos nx$, $x\in(0,1)$;
        \item $\displaystyle\sum\limits_{n=1}^{\infty}\dfrac{\ln(nx)}{1+n^3\ln^2x}$, $x\in\left(\dfrac{1}{2},1\right)$;
        \item $\displaystyle\sum\limits_{n=1}^{\infty}\dfrac{1}{(1+nx)^2}$, $x\in(0,1)$;
        \item $\displaystyle\sum\limits_{n=1}^{\infty}\left(1+n\dfrac{\sin x}{x}\right)^2$, $x\in(0,\pi)$.
    \end{enumerate}
\end{example}
\exampleblank

\begin{solution}
    \leavevmode
    \begin{enumerate}
        \item 当 $x\to0^+$ 时,各项趋于 $1$,相应和函数无界,故不可能是一致极限.
        \item 对每个 $x<1$ 级数收敛,但 $x\to1^-$ 时第 $n$ 项趋于 $\ln n$,端点处连通项趋零条件都失败,故不一致收敛.
        \item 当 $x\to0^+$ 时,$\displaystyle\sum(1+nx)^{-2}$ 发散到无穷;而各部分和有界,故不可能一致收敛.
        \item 当前通项对每个 $x\in(0,\pi)$ 都不趋于零,原级数逐点发散,不能称为``非一致收敛'';指数的正负号需要核对.
    \end{enumerate}
\end{solution}

%例题9.1.9
\begin{example}
    试用 Abel--Dirichlet 判别法判别下列级数在指定区间上的一致收敛性:
    \begin{enumerate}
        \item $\displaystyle\sum\limits_{n=1}^{\infty}\dfrac{(-1)^n(x^2+n)}{n^2}$, $x\in\lbrack a,b\rbrack$;
        \item $\displaystyle\sum\limits_{n=1}^{\infty}\dfrac{(-1)^{n+1}e^{-nx}}{\sqrt{n+x^2}}$, $x\in\lbrack 0,+\infty)$;
        \item $\displaystyle\sum\limits_{n=1}^{\infty}\dfrac{(-1)^n}{\sqrt n}\sin\left(1+\dfrac{x}{n}\right)$, $x\in\lbrack a,b\rbrack$;
        \item $\displaystyle\sum\limits_{n=1}^{\infty}\dfrac{\cos nx}{n^p}$, $x\in\lbrack \delta,2\pi-\delta\rbrack$, $p\in(0,1\rbrack$.
    \end{enumerate}
\end{example}
\exampleblank

\begin{solution}
    \leavevmode
    \begin{enumerate}
        \item 分解为 $x^2\displaystyle\sum(-1)^n/n^2+\sum(-1)^n/n$;前者由 M 判别,后者与 $x$ 无关,故一致收敛.
        \item $e^{-nx}/\sqrt{n+x^2}$ 关于 $n$ 递减,且其关于 $x\geq0$ 的上确界为 $1/\sqrt n\to0$,故一致收敛.
        \item 写成
              \[
                  \sin1\sum\limits_{n=1}^{\infty}\dfrac{(-1)^n}{\sqrt n}
                  +\sum\limits_{n=1}^{\infty}\dfrac{(-1)^n}{\sqrt n}
                  \left[\sin\left(1+\dfrac{x}{n}\right)-\sin1\right].
              \]
              第二级数在有限区间上由 $C/n^{3/2}$ 控制,故一致收敛.
        \item 三角和的部分和在 $[\delta,2\pi-\delta]$ 上一致有界,而 $n^{-p}\downarrow0$,由一致 Dirichlet 判别得结论.
    \end{enumerate}
\end{solution}

%例题9.1.10
\begin{example}
    试由和函数的不连续性判别下列级数在指定区间上的非一致收敛性:
    \begin{enumerate}
        \item $\displaystyle S(x)=\sum\limits_{n=0}^{\infty}\dfrac{x^4}{(1+x^4)^n}$, $x\in\lbrack 0,1\rbrack$;
        \item $\displaystyle S(x)=\sum\limits_{n=0}^{\infty}x^n(1-x)$, $x\in\lbrack 0,1\rbrack$;
        \item $\displaystyle S(x)=\sum\limits_{n=1}^{\infty}x^n\ln x$, $x\in(0,1\rbrack$;
        \item $\displaystyle S(x)=\sum\limits_{n=2}^{\infty}\left(x^{\frac{1}{2n-1}}-x^{\frac{1}{2n-3}}\right)$, $x\in\lbrack -1,1\rbrack$.
    \end{enumerate}
\end{example}
\exampleblank

\begin{solution}
    \leavevmode
    \begin{enumerate}
        \item 当 $x>0$ 时和为 $1+x^4$,而 $x=0$ 时和为 $0$,在零点不连续.
        \item 当 $0\leq x<1$ 时和为 $1$,而 $x=1$ 时和为 $0$,在右端点不连续.
        \item 当 $0<x<1$ 时和为 $x\ln x/(1-x)$,其在 $x\to1^-$ 时趋于 $-1$;但 $x=1$ 时每项为零,和为零.
        \item 前 $N$ 项裂项为 $x^{1/(2N-1)}-x$.极限在 $x>0$ 时为 $1-x$,在 $x=0$ 时为零,在 $x<0$ 时为 $-1-x$,故在零点不连续.
    \end{enumerate}
    各项均连续,而和函数不连续,所以均不一致收敛.
\end{solution}

%例题9.1.11
\begin{example}
    设 $f(x)$ 是 $(-\infty,+\infty)$ 上的连续函数,
    \[
        f_n(x)=\sum\limits_{k=0}^{n-1}\dfrac{1}{n}f\left(x+\dfrac{k}{n}\right).
    \]
    证明:函数列 $\{f_n(x)\}$ 在任何有限区间上一致收敛.
\end{example}
\exampleblank

\begin{solution}
    在任意有限区间 $[a,b]$ 上,$f$ 在 $[a,b+1]$ 上一致连续.$f_n(x)$ 是函数
    $t\mapsto f(x+t)$ 在 $[0,1]$ 上等距分割的 Riemann 和,故
    \[
        f_n(x)\rightrightarrows\int_0^1f(x+t)\,\mathrm{d}t
    \]
    且误差由 $f$ 在 $[a,b+1]$ 上的连续模一致控制.
\end{solution}

%例题9.1.12
\begin{example}
    设函数 $f(x)$ 在 $(-\infty,+\infty)$ 上有连续导函数 $f'(x)$,
    \[
        f_n(x)=e^n\left[f(x+e^{-n})-f(x)\right].
    \]
    证明:$\{f_n(x)\}$ 在任一有限开区间 $(a,b)$ 内一致收敛于 $f'(x)$.
\end{example}
\exampleblank

\begin{solution}
    写成
    \[
        f_n(x)=e^n\int_0^{e^{-n}}f'(x+t)\,\mathrm{d}t.
    \]
    在包含 $[a,b]$ 的稍大紧区间上,$f'$ 一致连续,故
    \[
        \sup_{a<x<b}|f_n(x)-f'(x)|
        \leq\sup_{\substack{a<x<b\\0\leq t\leq e^{-n}}}|f'(x+t)-f'(x)|\longrightarrow0.
    \]
\end{solution}

%例题9.1.13
\begin{example}[Heine 定理的推广]
    试证:$x\to a$ 时 $f(x,y)\rightrightarrows\varphi(y)$ 的充要条件是对任意 $\{x_n\}\to a$,有 $f(x_n,y)\rightrightarrows\varphi(y)$ $(n\to\infty)$.
\end{example}
\exampleblank

\begin{solution}
    必要性直接由定义得到.反之若不一致趋于 $\varphi$,则存在 $\varepsilon_0>0$,以及趋于 $a$ 的点列 $x_n$,使
    \[
        \sup_y|f(x_n,y)-\varphi(y)|\geq\varepsilon_0.
    \]
    这与题设对任意点列均一致收敛矛盾,故条件充分.
\end{solution}

%例题9.1.14
\begin{example}
    若 $f_n(x)$ 在 $\lbrack a,b\rbrack$ 上可积,$f(x)$ 与 $g(x)$ 在 $\lbrack a,b\rbrack$ 上都可积,且
    \[
        \displaystyle\lim\limits_{n\to\infty}\int_a^b|f_n(x)-f(x)|^2\,\mathrm{d}x=0.
    \]
    设
    \[
        h(x)=\int_a^x f(t)g(t)\,\mathrm{d}t,
        \qquad
        h_n(x)=\int_a^x f_n(t)g(t)\,\mathrm{d}t.
    \]
    证明:$h_n(x)$ 在 $\lbrack a,b\rbrack$ 上一致收敛于 $h(x)$.
\end{example}
\exampleblank

\begin{solution}
    当前条件不足.取区间 $[0,1]$,令 $f=0$,$g(x)=x^{-2/3}$,并令
    \[
        f_n(x)=n^{1/3}\mathbf1_{(0,1/n)}(x).
    \]
    则 $g$ 可积,且
    \[
        \displaystyle\int_0^1|f_n|^2=n^{-1/3}\to0,
    \]
    但 $h_n(1)=\displaystyle\int_0^1f_ng=3$,并不趋于零.若补充 $g\in L^2[a,b]$,则由 Cauchy--Schwarz 不等式
    \[
        \sup_x|h_n(x)-h(x)|
        \leq\|f_n-f\|_2\|g\|_2\longrightarrow0.
    \]
    故题面需要核对.
\end{solution}

%例题9.1.15
\begin{example}
    设 $f_1(x)$ 在 $\lbrack a,b\rbrack$ 上正常可积,
    \[
        f_{n+1}(x)=\int_a^x f_n(t)\,\mathrm{d}t,
        \qquad n=1,2,\ldots.
    \]
    证明:函数序列 $\{f_n(x)\}$ 在 $\lbrack a,b\rbrack$ 上一致收敛于零.
\end{example}
\exampleblank

\begin{solution}
    设 $M=\sup_{[a,b]}|f_1|$,$L=b-a$.反复积分可得
    \[
        |f_n(x)|\leq M\dfrac{L^{n-1}}{(n-1)!},
        \qquad x\in[a,b].
    \]
    右端趋于零且与 $x$ 无关,故 $f_n$ 一致收敛于零.
\end{solution}

%例题9.1.16
\begin{example}
    \begin{enumerate}
        \item 设 $f_n(x)$ 在 $\lbrack a,b\rbrack$ 上可微,且存在 $M>0$,使 $|f_n'(x)|\leq M$.若 $f_n(x)$ 在 $\lbrack a,b\rbrack$ 上逐点收敛于 $f(x)$,则必一致收敛;
        \item 设定义在 $\lbrack 0,1\rbrack$ 上的函数列 $\{f_n(x)\}$ 满足存在 $M>0$ 使得
              $|f_n(x)-f_n(y)|\leq M|x-y|$,且 $\displaystyle\lim\limits_{n\to\infty}f_n(x)=f(x)$,则 $f_n(x)$ 在 $\lbrack 0,1\rbrack$ 上一致收敛于 $f(x)$.
    \end{enumerate}
\end{example}
\exampleblank

\begin{solution}
    两题本质相同.由统一的 Lipschitz 估计,函数族等度连续;逐点极限 $f$ 也满足同一 Lipschitz 估计.任取 $\varepsilon>0$,先把区间分成有限个长度小于 $\varepsilon/(3M)$ 的小区间,再在所有分点处利用逐点收敛统一选取 $N$.任意点与相邻分点比较,即得 $n\geq N$ 时
    \[
        \sup_x|f_n(x)-f(x)|<\varepsilon.
    \]
    故两列均一致收敛.
\end{solution}

%例题9.1.17
\begin{example}
    设 $\{f_n(x)\}$ 是 $\lbrack 0,1\rbrack$ 上由递增函数组成的序列,若存在 $f\in C\lbrack 0,1\rbrack$ 使得 $\displaystyle\lim\limits_{n\to\infty}f_n(x)=f(x)$,则 $f_n(x)$ 在 $\lbrack 0,1\rbrack$ 上一致收敛于 $f(x)$.
\end{example}
\exampleblank

\begin{solution}
    任取 $\varepsilon>0$.由 $f$ 的一致连续性,取分点
    $0=x_0<\cdots<x_m=1$,使相邻分点处 $f$ 的振幅小于 $\varepsilon$.在有限个分点上同时取充分大的 $n$,使
    $|f_n(x_j)-f(x_j)|<\varepsilon$.若 $x_j\leq x\leq x_{j+1}$,由
    $f_n$ 的递增性,
    \[
        f_n(x_j)\leq f_n(x)\leq f_n(x_{j+1}),
    \]
    结合 $f$ 的振幅估计得到 $|f_n(x)-f(x)|<2\varepsilon$,故一致收敛.
\end{solution}

%例题9.1.18
\begin{example}
    设 $\{a_n\}$ 是单调收敛于 $0$ 的数列,则
    \begin{enumerate}
        \item $\displaystyle\sum\limits_{n=1}^{\infty}(-1)^{n-1}a_n\cos nx$ 在 $\lbrack -\pi+\delta,\pi-\delta\rbrack$ $(\delta>0)$ 上一致收敛;
        \item $\displaystyle\sum\limits_{n=1}^{\infty}a_nx\sin nx$ 在 $\lbrack -\pi+\delta,\pi-\delta\rbrack$ $(\delta>0)$ 上一致收敛.
    \end{enumerate}
\end{example}
\exampleblank

\begin{solution}
    \leavevmode
    \begin{enumerate}
        \item $(-1)^{n-1}\cos(nx)$ 的部分和在给定闭区间上一致有界,故由一致 Dirichlet 判别法收敛.
        \item 由三角和公式,
              \[
                  \left|x\sum\limits_{k=1}^{N}\sin(kx)\right|
                  \leq C_\delta
              \]
              在给定区间上一致成立;再与 $a_n\downarrow0$ 配合使用 Abel 变换,即得一致收敛.
    \end{enumerate}
\end{solution}

%例题9.1.19
\begin{example}
    \begin{enumerate}
        \item 若 $\{na_n\}$ 是单调收敛于 $0$ 的数列,则 $\displaystyle\sum\limits_{n=1}^{\infty}a_n\sin nx$ 在 $(-\infty,+\infty)$ 上一致收敛;
        \item 设 $\{a_n\}$ 是递减正数列.若 $\displaystyle\sum\limits_{n=1}^{\infty}a_n\sin nx$ 在 $\lbrack -a,a\rbrack$ 上一致收敛,则 $na_n\to0$ $(n\to\infty)$.
    \end{enumerate}
\end{example}
\exampleblank

\begin{solution}
    \leavevmode
    \begin{enumerate}
        \item 这是 Chaundy--Jolliffe 判别.令 $m=[1/|x|]$:在 $n\leq m$ 的部分使用 $|\sin(nx)|\leq n|x|$,在 $n>m$ 的部分使用 Abel 变换与三角和估计;两部分均由 $\sup_{k\geq N}k|a_k|$ 控制,故尾和在全实轴上一致趋零.
        \item 由一致 Cauchy 准则,取 $x=1/n$ 并考察 $n\leq k\leq2n$,有
              \[
                  \displaystyle\sum\limits_{k=n}^{2n}a_k\sin\dfrac{k}{n}\geq cna_{2n}\longrightarrow0.
              \]
              所以 $na_{2n}\to0$;再由单调性夹逼得到 $na_n\to0$.
    \end{enumerate}
\end{solution}

%例题9.1.20
\begin{example}
    设 $\displaystyle\sum\limits_{n=1}^{\infty}a_n$ 收敛,且 $\{f_n(x)\}$ 满足
    \[
        M\geq f_1(x)\geq f_2(x)\geq\cdots\geq f_n(x)\geq\cdots\geq0.
    \]
    证明:$\displaystyle\sum\limits_{n=1}^{\infty}a_nf_n(x)$ 在 $\lbrack a,b\rbrack$ 上一致收敛.
\end{example}
\exampleblank

\begin{solution}
    记 $A_n=\displaystyle\sum_{k=1}^{n}a_k$.对任意 $m>n$ 作 Abel 变换,尾和可写成
    \[
        \displaystyle\sum\limits_{k=n}^{m}a_kf_k
        =(A_m-A_{n-1})f_m+
        \sum\limits_{k=n}^{m-1}(A_k-A_{n-1})(f_k-f_{k+1}).
    \]
    由于 $\displaystyle\sum a_n$ 的尾部分和一致小,而 $f_k$ 的总变差不超过 $M$,右端一致趋于零,故级数一致收敛.
\end{solution}

%例题9.1.21
\begin{example}
    设定义在区间 $I$ 上的 $\{f_n(x)\},\{g_n(x)\}$ 满足:
    \begin{enumerate}
        \item $\displaystyle\sum\limits_{n=1}^{\infty}|f_{n+1}(x)-f_n(x)|$ 在 $I$ 上一致收敛;
        \item $\displaystyle\lim\limits_{n\to\infty}\sup\limits_I|f_n(x)|=0$;
        \item $G_n=\displaystyle\sum\limits_{k=1}^n g_k(x)$ 在 $I$ 上一致有界.
    \end{enumerate}
    则 $\displaystyle\sum\limits_{n=1}^{\infty}f_n(x)g_n(x)$ 在 $I$ 上一致收敛.
\end{example}
\exampleblank

\begin{solution}
    对尾和作分部求和:
    \[
        \displaystyle\sum\limits_{k=n}^{m}f_kg_k
        =f_mG_m-f_nG_{n-1}
        +\sum\limits_{k=n}^{m-1}(f_k-f_{k+1})G_k.
    \]
    由 $G_k$ 一致有界,$f_n\rightrightarrows0$,以及
    $\displaystyle\sum|f_{n+1}-f_n|$ 的一致收敛性,右端对 $m>n$ 一致趋零.故结论由一致 Cauchy 准则成立.
\end{solution}

\begin{definition}
    设 $\{f_n(x)\}$ 是定义在区间 $I$ 上的函数列.若对任给 $\varepsilon>0$,存在 $\delta>0$,使得当 $I$ 中的点 $x',x''$ 满足 $|x'-x''|<\delta$ 时,对一切正整数 $n$,都有 $|f_n(x')-f_n(x'')|<\varepsilon$,则称 $\{f_n(x)\}$ 在 $I$ 上是等度连续的.
\end{definition}

\begin{remark}
    易知等度连续函数列中的每一个函数都在 $I$ 上一致连续.
\end{remark}

\begin{theorem}[
        \markstar]
    设 $\{f_n(x)\}$ 为 $I$ 上函数列,$E$ 在 $I$ 中稠密且 $\{f_n(x)\}$ 在 $E$ 上逐点收敛,则存在子列 $\{f_{n_k}(x)\}$ 在 $E$ 上逐点收敛.
\end{theorem}

\begin{theorem}[\markstar]
    \leavevmode
    \begin{enumerate}
        \item 若 $\lbrack a,b\rbrack$ 上连续函数列 $\{f_n(x)\}$ 一致收敛,则 $\{f_n(x)\}$ 等度连续;
        \item 若 $\lbrack a,b\rbrack$ 上等度连续函数列 $\{f_n(x)\}$ 满足 $f_n(x)\to f(x)$ $(n\to\infty)$,则 $f_n(x)\rightrightarrows f(x)$ $(n\to\infty)$.
    \end{enumerate}
\end{theorem}

\begin{theorem}[\markstar]
    设 $\lbrack a,b\rbrack$ 上等度连续函数列 $\{f_n(x)\}$ 在 $\lbrack a,b\rbrack$ 逐点有界,则
    \begin{enumerate}
        \item $\{f_n(x)\}$ 在 $\lbrack a,b\rbrack$ 上一致有界;
        \item $\{f_n(x)\}$ 有一致收敛的子列 $\{f_{n_k}(x)\}$.
    \end{enumerate}
\end{theorem}

%例题9.1.22
\begin{example}
    设 $\{f_n(x)\}$ 是区间 $(a,b)$ 内的连续函数序列,并且对任意 $x_0\in(a,b)$,$\{f_n(x_0)\}$ 都是有界的.证明:$\{f_n(x)\}$ 在 $(a,b)$ 的某一非空子区间上一致有界.
\end{example}
\exampleblank

\begin{solution}
    令
    \[
        E_m=\{x\in(a,b):|f_n(x)|\leq m\text{ 对所有 }n\}.
    \]
    各 $E_m$ 为相对闭集,且由逐点有界性有 $(a,b)=\displaystyle\bigcup_{m=1}^{\infty}E_m$.由 Baire 类定理,某个 $E_m$ 在 $(a,b)$ 中有非空内点;在该开子区间上,全体 $f_n$ 一致有界于 $m$.
\end{solution}

%例题9.1.23
\begin{example}
    设 $\{f_n(x)\}$ 在区间 $\lbrack 0,1\rbrack$ 上一致有界,试证存在一子列,其在 $\lbrack 0,1\rbrack$ 的一切有理点上收敛.
\end{example}
\exampleblank

\begin{solution}
    将 $[0,1]$ 内有理数排列为 $r_1,r_2,\ldots$.由一致有界性,先选子列使其在 $r_1$ 处收敛,再从中选子列使其在 $r_2$ 处收敛,如此继续.取对角子列 $f_{n_k}$,则对每个固定 $j$,它从某项起属于第 $j$ 次所选子列,故在 $r_j$ 处收敛.
\end{solution}

%例题9.1.24
\begin{example}
    设 $\mathscr R$ 是区间 $I$ 上的函数族,对任意 $f\in\mathscr R$ 都在 $I$ 上可微,且 $\{f'(x)\mid f\in\mathscr R\}$ 在 $I$ 上一致有界.证明:$\mathscr R$ 在 $I$ 上等度连续.
\end{example}
\exampleblank

\begin{solution}
    若 $|f'(x)|\leq M$ 对所有 $f\in\mathscr R$,$x\in I$ 成立,则由 Lagrange 中值定理
    \[
        |f(x)-f(y)|\leq M|x-y|.
    \]
    给定 $\varepsilon>0$,统一取 $\delta=\varepsilon/M$ 即得等度连续性.
\end{solution}

%例题9.1.25
\begin{example}
    \begin{enumerate}
        \item 设 $\{f_n(x)\}$ 是 $\lbrack 0,1\rbrack$ 上二次可导的函数列,且 $f_n(0)=f_n'(0)=0$.若 $|f_n''(x)|\leq1$,则存在子列 $\{f_{n_k}(x)\}$ 在 $\lbrack 0,1\rbrack$ 上一致收敛;
        \item 设 $f_n\in C^{(1)}\lbrack 0,1\rbrack$,且 $|f_n'(x)|\leq x^{-\frac{1}{2}}$ $(0<x\leq1)$,并且 $\displaystyle\int_0^1f_n(x)\,\mathrm{d}x=0$,则存在子列 $\{f_{n_k}(x)\}$ 在 $\lbrack 0,1\rbrack$ 上一致收敛.
    \end{enumerate}
\end{example}
\exampleblank

\begin{solution}
    \leavevmode
    \begin{enumerate}
        \item 由初值与二阶导数界,$|f_n'(x)|\leq1$ 且 $|f_n(x)|\leq1/2$.故函数族一致有界且等度连续,由 Arzel\`a--Ascoli 定理存在一致收敛子列.
        \item 对 $0\leq x<y\leq1$,
              \[
                  |f_n(y)-f_n(x)|\leq\int_x^yt^{-1/2}\,\mathrm{d}t
                  \leq2\sqrt{y-x},
              \]
              故等度连续.又由积分为零,存在 $c_n\in[0,1]$ 使 $f_n(c_n)=0$,于是 $|f_n(x)|\leq2$.再次应用 Arzel\`a--Ascoli 定理即得一致收敛子列.
    \end{enumerate}
\end{solution}


\newpage

\section{函数性质的传递----极限换序}


\subsection{连续性质的传递}

\begin{theorem}[\markstar]
    若 $u_n(x)$ 在区间 $I$ 上连续,$\displaystyle\sum\limits_{n=1}^{\infty}u_n(x)$ 在 $I$ 上一致收敛,则
    \[
        S(x)=\sum\limits_{n=1}^{\infty}u_n(x)
    \]
    在 $I$ 上连续.
\end{theorem}

\begin{theorem}[\markstar]
    若 $u_n(x)$ 在 $x=x_0$ 处连续,$\displaystyle\sum\limits_{n=1}^{\infty}u_n(x)$ 在 $x_0$ 的某个邻域里一致收敛,则 $S(x)=\displaystyle\sum\limits_{n=1}^{\infty}u_n(x)$ 在 $x=x_0$ 处连续.
\end{theorem}

\begin{theorem}[\markstar]
    若 $u_n(x)$ 在 $(a,b)$ 内连续,$\displaystyle\sum\limits_{n=1}^{\infty}u_n(x)$ 在 $(a,b)$ 内闭一致收敛,则 $S(x)=\displaystyle\sum\limits_{n=1}^{\infty}u_n(x)$ 在 $(a,b)$ 内连续.
\end{theorem}

\begin{remark}
    一致收敛是和函数连续的充分不必要条件.例如 $\lbrack 0,1\rbrack$ 上的 $f_n(x)=nxe^{-nx^2}$ 收敛到 $f(x)=0$,但非一致收敛.
\end{remark}

%例题9.2.1
\begin{example}
    判别下列级数 $\displaystyle S(x)=\sum\limits_{n=1}^{\infty}f_n(x)$ 在指定区间上的连续性:
    \begin{enumerate}
        \item $\displaystyle\sum\limits_{n=1}^{\infty}\dfrac{\ln^n x}{n^2}$, $x\in\lbrack e^{-1},e\rbrack$;
        \item $\displaystyle\sum\limits_{n=1}^{\infty}\dfrac{x+n(-1)^n}{n^2+x^2}$, $x\in\lbrack -a,a\rbrack$;
        \item $\displaystyle\sum\limits_{n=1}^{\infty}\dfrac{x^n\sin nx}{n!}$, $x\in\lbrack a,b\rbrack$;
        \item $\displaystyle\sum\limits_{n=1}^{\infty}\left(\dfrac{1}{n}+x\right)^n$, $x\in(-1,1)$.
    \end{enumerate}
\end{example}
\exampleblank

\begin{solution}
    \leavevmode
    \begin{enumerate}
        \item 因 $|\ln x|\leq1$,通项绝对值由 $1/n^2$ 控制,故和函数连续.
        \item 分成
              \[
                  \displaystyle\sum\limits_{n=1}^{\infty}\dfrac{x}{n^2+x^2}
                  +\sum\limits_{n=1}^{\infty}(-1)^n\dfrac{n}{n^2+x^2}.
              \]
              第一项在紧区间上一致绝对收敛;第二项从 $n>a$ 起系数关于 $n$ 递减且上确界趋零,故一致收敛.和函数连续.
        \item 通项由 $C^n/n!$ 控制,故在任意有限闭区间上一致收敛,和函数连续.
        \item 对任意 $x_0\in(-1,1)$,取 $|x_0|<r<1$,则在 $x_0$ 的充分小邻域内,充分大时 $|x+1/n|\leq q<1$.故级数局部一致收敛,和函数在 $(-1,1)$ 上连续.
    \end{enumerate}
\end{solution}

%例题9.2.2
\begin{example}
    试求下列极限值:
    \begin{enumerate}
        \item $\displaystyle\lim\limits_{x\to0^+}\sum\limits_{n=1}^{\infty}\dfrac{1}{2^n n^x}$;
        \item $\displaystyle\lim\limits_{x\to+\infty}\sum\limits_{n=1}^{\infty}\dfrac{x^2}{1+n^2x^2}$;
        \item $\displaystyle\lim\limits_{x\to1^-}\sum\limits_{n=1}^{\infty}\dfrac{(-1)^{n+1}}{n}\dfrac{x^n}{1+x^n}$;
        \item $\displaystyle\lim\limits_{x\to1}\sum\limits_{n=1}^{\infty}\dfrac{x^n\sin\left(\dfrac{n\pi x}{2}\right)}{2^n}$.
    \end{enumerate}
\end{example}
\exampleblank

\begin{solution}
    利用一致收敛或控制收敛逐项取极限:
    \begin{enumerate}
        \item 由 $2^{-n}$ 控制,极限为 $\displaystyle\sum_{n=1}^{\infty}2^{-n}=\boxed1$.
        \item 因 $0\leq x^2/(1+n^2x^2)\leq1/n^2$,极限为
              $\displaystyle\sum_{n=1}^{\infty}1/n^2=\boxed{\pi^2/6}$.
        \item 一致 Abel 判别允许令 $x\to1^-$,故极限为
              \[
                  \dfrac{1}{2}\sum\limits_{n=1}^{\infty}\dfrac{(-1)^{n+1}}n
                  =\boxed{\dfrac{1}{2}\ln2}.
              \]
        \item 由几何级数控制,令 $x\to1$ 得
              \[
                  \displaystyle\sum\limits_{n=1}^{\infty}\dfrac{\sin(n\pi/2)}{2^n}
                  =\operatorname{Im}\dfrac{i/2}{1-i/2}
                  =\boxed{\dfrac{2}{5}}.
              \]
    \end{enumerate}
\end{solution}

%例题9.2.3
\begin{example}
    设 $g_n,h_n\in C\lbrack a,b\rbrack$,且 $g_n(x)\leq g_{n+1}(x)$,$h_n(x)\geq h_{n+1}(x)$,
    \[
        \displaystyle\lim\limits_{n\to\infty}g_n(x)=f(x)=\lim\limits_{n\to\infty}h_n(x).
    \]
    则 $f\in C\lbrack a,b\rbrack$.
\end{example}
\exampleblank

\begin{solution}
    令 $d_n=h_n-g_n\geq0$.则 $d_n$ 连续,关于 $n$ 单调递减,并逐点趋于零.由 Dini 定理,$d_n\rightrightarrows0$.于是
    \[
        0\leq f-g_n\leq d_n,
    \]
    故 $g_n\rightrightarrows f$.连续函数的一致极限仍连续,因而 $f\in C[a,b]$.
\end{solution}

%例题9.2.4
\begin{example}
    设 $f_n,F_n\in C\lbrack a,b\rbrack$ 且 $|f_n(x)|\leq F_n(x)$.若
    \[
        \sigma(x)=\sum\limits_{n=1}^{\infty}F_n(x)
    \]
    在 $\lbrack a,b\rbrack$ 上连续,则 $S(x)=\displaystyle\sum\limits_{n=1}^{\infty}f_n(x)$ 在 $\lbrack a,b\rbrack$ 上连续.
\end{example}
\exampleblank

\begin{solution}
    由 $|f_n|\leq F_n$ 可知 $F_n\geq0$.连续函数的部分和
    $\sigma_N=\displaystyle\sum_{n=1}^{N}F_n$ 单调增加到连续函数 $\sigma$,故由 Dini 定理一致收敛.因此
    \[
        \sup_{x\in[a,b]}\sum\limits_{n=N+1}^{\infty}|f_n(x)|
        \leq\sup_{x\in[a,b]}\left(\sigma(x)-\sigma_N(x)\right)\longrightarrow0.
    \]
    所以 $\displaystyle\sum f_n$ 一致收敛,其和函数连续.
\end{solution}

%例题9.2.5
\begin{example}
    证明:$\displaystyle\sum\limits_{n=-\infty}^{+\infty}\dfrac{1}{(n-x)^2}$ 当 $x$ 不为整数时收敛,周期为 $1$,并且当 $x$ 非整数时和函数连续.
\end{example}
\exampleblank

\begin{solution}
    在任一不含整数的紧区间 $K$ 上,除有限项外有
    \[
        \dfrac{1}{(n-x)^2}\leq\dfrac{4}{n^2}\qquad(x\in K),
    \]
    故双向级数在 $K$ 上一致收敛,和函数在每个非整数点连续.换指标
    $m=n-1$ 得
    \[
        S(x+1)=\sum\limits_{n\in\mathbb{Z}}\dfrac{1}{(n-x-1)^2}
        =\sum\limits_{m\in\mathbb{Z}}\dfrac{1}{(m-x)^2}=S(x),
    \]
    故周期为 $1$.
\end{solution}

%例题9.2.6
\begin{example}[\markstar]
    设 $\{x_n\}$ 是 $(0,1)$ 内的一个序列:$0<x_n<1$ 且 $x_i\neq x_j$ $(i\neq j)$.试讨论函数
    \[
        f(x)=\sum\limits_{n=1}^{\infty}\dfrac{\operatorname{sgn}(x-x_n)}{2^n}
    \]
    在 $(0,1)$ 内的连续性.
\end{example}
\exampleblank

\begin{solution}
    级数由 $2^{-n}$ 控制而一致收敛.固定 $x\notin\{x_n\}$,给定
    $\varepsilon>0$,先使尾和小于 $\varepsilon$,再取一个不含有限个首项跳点的邻域,可知 $f$ 在 $x$ 连续.

    在 $x=x_k$ 处,第 $k$ 项的左右极限相差 $2/2^k=2^{1-k}$;其余有限个首项在该点附近不变,尾项可一致控制.因此 $f$ 在每个 $x_k$ 处有非零跳跃并不连续.故不连续点集恰为 $\{x_n:n\geq1\}$.
\end{solution}



\subsection{积分性质的传递}

一致收敛性只是积分换序的充分不必要条件.换序的优越性在于不必求出和函数而可以计算它的积分值.进一步,
\[
    \begin{aligned}
        \int_a^b\sum\limits_{k=1}^{\infty}u_k(x)\,\mathrm{d}x
         & =\sum\limits_{k=1}^{\infty}\int_a^b u_k(x)\,\mathrm{d}x                                                                                  \\
         & \iff \lim\limits_{n\to\infty}\left|\int_a^b\left(\sum\limits_{k=1}^{\infty}u_k(x)-\sum\limits_{k=1}^{n}u_k(x)\right)\mathrm{d}x\right|=0 \\
         & \iff \lim\limits_{n\to\infty}\int_a^bR_n(x)\,\mathrm{d}x=0.
    \end{aligned}
\]

%例题9.2.7
\begin{example}
    设 $f\in C^{\infty}(-\infty,+\infty)$,函数列 $\{f^{(n)}(x)\}$ 在任意区间 $\lbrack a,b\rbrack$ 上一致收敛于 $g(x)$,则 $g(x)=Ce^x$.
\end{example}
\exampleblank

\begin{solution}
    对任意 $a<x<b$,
    \[
        f^{(n)}(x)-f^{(n)}(a)=\int_a^xf^{(n+1)}(t)\,\mathrm{d}t.
    \]
    两列 $f^{(n)}$ 与 $f^{(n+1)}$ 都一致收敛于 $g$,故令 $n\to\infty$ 得
    \[
        g(x)-g(a)=\int_a^xg(t)\,\mathrm{d}t.
    \]
    于是 $g'=g$,故 $g(x)=Ce^x$.
\end{solution}

%例题9.2.8
\begin{example}
    证明下列命题:
    \begin{enumerate}
        \item $\displaystyle\int_0^1\sum\limits_{n=1}^{\infty}\dfrac{x}{n(n+x)}\,\mathrm{d}x=\gamma$,其中 $\gamma$ 为 Euler 常数;
        \item $\displaystyle\int_{\ln2}^{\ln5}\sum\limits_{n=1}^{\infty}ne^{-nx}\,\mathrm{d}x=\dfrac{3}{4}$;
        \item
              \[
                  \displaystyle\int_0^{+\infty}\dfrac{x^{a-1}}{1+x}\,\mathrm{d}x
                  =\dfrac{1}{a}+\sum\limits_{k=1}^{\infty}(-1)^k
                  \left(\dfrac{1}{a+k}-\dfrac{1}{a-k}\right),
                  \qquad 0<a<1.
              \]
    \end{enumerate}
\end{example}
\exampleblank

\begin{solution}
    \leavevmode
    \begin{enumerate}
        \item 因
              \[
                  \dfrac{x}{n(n+x)}=\dfrac{1}{n}-\dfrac{1}{n+x},
              \]
              逐项积分得到
              \[
                  \displaystyle\sum\limits_{n=1}^{\infty}\left[\dfrac{1}{n}-\ln\left(1+\dfrac{1}{n}\right)\right]=\gamma.
              \]
        \item 逐项积分得
              \[
                  \displaystyle\sum\limits_{n=1}^{\infty}(2^{-n}-5^{-n})=1-\dfrac{1}{4}=\dfrac{3}{4}.
              \]
        \item 分别在 $(0,1)$ 上展开 $1/(1+x)$,并对 $(1,+\infty)$ 作代换 $t=1/x$ 后展开.逐项积分并合并,得到
              \[
                  \dfrac{1}{a}+\sum\limits_{k=1}^{\infty}(-1)^k
                  \left(\dfrac{1}{a+k}-\dfrac{1}{a-k}\right).
              \]
              在 $0<a<1$ 时上述换序可先在 $[\varepsilon,1-\varepsilon]$ 上进行,再令 $\varepsilon\to0^+$.
    \end{enumerate}
\end{solution}

%例题9.2.9
\begin{example}
    证明下列命题:
    \begin{enumerate}
        \item $\displaystyle\lim\limits_{n\to\infty}\int_0^1\dfrac{\mathrm{d}x}{1+\left(\dfrac{x}{n}+1\right)^n}=\ln\dfrac{2e}{e+1}$;
        \item $\displaystyle\int_0^1x^{-x}\,\mathrm{d}x=\sum\limits_{n=1}^{\infty}\dfrac{1}{n^n}$.
    \end{enumerate}
\end{example}
\exampleblank

\begin{solution}
    \leavevmode
    \begin{enumerate}
        \item 被积函数逐点趋于 $1/(1+e^x)$ 且由 $1$ 控制,故
              \[
                  \displaystyle\lim\limits_{n\to\infty}\int_0^1\dfrac{\mathrm{d}x}{1+(1+x/n)^n}
                  =\int_0^1\dfrac{\mathrm{d}x}{1+e^x}
                  =\ln\dfrac{2e}{e+1}.
              \]
        \item 展开
              \[
                  x^{-x}=e^{-x\ln x}
                  =\sum\limits_{k=0}^{\infty}\dfrac{x^k(-\ln x)^k}{k!}.
              \]
              各项非负,可逐项积分;又
              \[
                  \displaystyle\int_0^1x^k(-\ln x)^k\,\mathrm{d}x=\dfrac{k!}{(k+1)^{k+1}},
              \]
              故积分等于 $\displaystyle\sum_{n=1}^{\infty}n^{-n}$.
    \end{enumerate}
\end{solution}

%例题9.2.10
\begin{example}
    设 $f,g\in C(-\infty,+\infty)$,且 $f(x+1)=f(x)$,$g(x+1)=g(x)$.证明:
    \[
        \displaystyle\lim\limits_{n\to\infty}\int_0^1f(x)g(nx)\,\mathrm{d}x
        =\left(\int_0^1f(x)\,\mathrm{d}x\right)
        \left(\int_0^1g(x)\,\mathrm{d}x\right).
    \]
\end{example}
\exampleblank

\begin{solution}
    把 $[0,1]$ 分成 $n$ 个小区间并令 $x=(k+t)/n$,得到
    \[
        \displaystyle\int_0^1f(x)g(nx)\,\mathrm{d}x
        =\int_0^1g(t)\left[\dfrac{1}{n}\sum\limits_{k=0}^{n-1}f\left(\dfrac{k+t}{n}\right)\right]\mathrm{d}t.
    \]
    方括号内的 Riemann 和关于 $t\in[0,1]$ 一致趋于 $\displaystyle\int_0^1f$,故可取极限,得到题中乘积.
\end{solution}

%例题9.2.11
\begin{example}
    设 $h(x),f_n'(x)$ 在 $\lbrack a,b\rbrack$ 上连续,又对 $\lbrack a,b\rbrack$ 中任意 $x_1,x_2$ 和正整数 $n$,有
    \[
        |f_n(x_1)-f_n(x_2)|\leq\dfrac{M}{n}|x_1-x_2|,
        \qquad M>0.
    \]
    证明:
    \[
        \displaystyle\lim\limits_{n\to\infty}\int_a^bh(x)f_n'(x)\,\mathrm{d}x=0.
    \]
\end{example}
\exampleblank

\begin{solution}
    由给定 Lipschitz 条件及 $f_n'$ 连续,令 $x_2\to x_1$ 得
    \[
        |f_n'(x)|\leq\dfrac{M}{n}.
    \]
    因此
    \[
        \left|\int_a^bh(x)f_n'(x)\,\mathrm{d}x\right|
        \leq\dfrac{M(b-a)}n\max_{[a,b]}|h|\longrightarrow0.
    \]
\end{solution}

%例题9.2.12
\begin{example}
    设 $g(x),f_n(x)$ 在 $\lbrack a,b\rbrack$ 上有界可积,且对任意 $c\in(a,b)$,$f_n(x)\rightrightarrows0$ 于 $\lbrack c,b\rbrack$ 上,
    \[
        \displaystyle\lim\limits_{n\to\infty}\int_a^bf_n(x)\,\mathrm{d}x=1,
        \qquad
        \lim\limits_{x\to a^+}g(x)=A.
    \]
    试证:
    \[
        \displaystyle\lim\limits_{n\to\infty}\int_a^bg(x)f_n(x)\,\mathrm{d}x=A.
    \]
\end{example}
\exampleblank

\begin{solution}
    当前条件不足,因为没有控制 $\displaystyle\int|f_n|$.例如在 $[0,1]$ 上取 $g(x)=x$,可在 $(0,1/n)$ 内取两个符号相反,总积分为零而一阶矩恒为非零的尖峰,再加上积分为 $1$ 的正尖峰;此时仍满足远离零一致趋零及 $\displaystyle\int f_n=1$,但 $\displaystyle\int gf_n$ 不趋于 $g(0)$.

    若补充 $f_n\geq0$(或 $\displaystyle\int|f_n|$ 一致有界),则任取 $c>a$,把积分分为 $[a,c]$ 与 $[c,b]$.前一段用 $g(x)\to A$ 控制,后一段用 $f_n\rightrightarrows0$ 控制,再令 $c\to a^+$,即可得到极限 $A$.故原题需要补充条件.
\end{solution}

%例题9.2.13
\begin{example}
    设 $\{f_n(x)\}$ 是定义在 $\lbrack -1,1\rbrack$ 上的连续函数序列,在原点附近一致有界.若对任意 $0<\delta<1$,$f_n(x)$ 在 $\lbrack -1,-\delta\rbrack\cup\lbrack \delta,1\rbrack$ 上一致收敛于零,且 $g\in C\lbrack -1,1\rbrack$, $g(0)=0$,证明:
    \[
        \displaystyle\lim\limits_{n\to\infty}\int_{-1}^{1}f_n(x)g(x)\,\mathrm{d}x=g(0).
    \]
\end{example}
\exampleblank

\begin{solution}
    结论右端 $g(0)=0$.取 $\delta>0$.在 $|x|<\delta$ 上,$f_n$ 一致有界且
    $\sup_{|x|<\delta}|g(x)|\to0$($\delta\to0$);在其余闭集上,$f_n$ 一致趋零.因此
    \[
        \left|\int_{-1}^{1}f_n(x)g(x)\,\mathrm{d}x\right|
        \leq C\int_{-\delta}^{\delta}|g(x)|\,\mathrm{d}x+o(1).
    \]
    先令 $n\to\infty$,再令 $\delta\to0$,即得极限为零.
\end{solution}



\subsection{微分性质的传递}

一致收敛对于微分性质的传递只是充分不必要条件.例如 $\displaystyle f(x)=\sum\limits_{n=1}^{\infty}ne^{-nx}$ 在 $(0,+\infty)$ 上非一致收敛,但和函数在其上无穷次可微.

%例题9.2.14
\begin{example}
    讨论下列级数 $\displaystyle S(x)=\sum\limits_{n=1}^{\infty}f_n(x)$ 在指定区域的可微性:
    \begin{enumerate}
        \item $f_n(x)=\arctan\dfrac{x}{n^2}$, $x\in(0,+\infty)$;
        \item $f_n(x)=\dfrac{\sin nx}{n^3}$, $x\in(-\infty,+\infty)$;
        \item $f_n(x)=\sqrt n\tan^n x$, $x\in\left(-\dfrac{\pi}{4},\dfrac{\pi}{4}\right)$.
    \end{enumerate}
\end{example}
\exampleblank

\begin{solution}
    \leavevmode
    \begin{enumerate}
        \item
              \[
                  f_n'(x)=\dfrac{n^{-2}}{1+x^2/n^4},\qquad |f_n'(x)|\leq n^{-2}.
              \]
              导数级数在全区间一致收敛,故和函数可微.
        \item $f_n'(x)=\cos(nx)/n^2$,导数级数由 $\displaystyle\sum n^{-2}$ 一致控制,故和函数在实轴上可微.
        \item 在任一紧子区间 $|\tan x|\leq q<1$ 上,原级数及导数级数都由多项式因子乘 $q^n$ 控制,故局部一致收敛并可逐项求导.因此和函数在给定开区间内可微.
    \end{enumerate}
\end{solution}

%例题9.2.15
\begin{example}
    设 $f(x)$ 在 $(-\infty,+\infty)$ 内有任意阶导数,级数
    \[
        \begin{aligned}
            \cdots+f^{(n)}(x)+\cdots+f''(x)+f'(x)+f(x)
             & +\int_0^xf(t_1)\,\mathrm{d}t_1                                 \\
             & +\int_0^x\mathrm{d}t_2\int_0^{t_2}f(t_1)\,\mathrm{d}t_1+\cdots
        \end{aligned}
    \]
    按两个方向在 $(-\infty,+\infty)$ 内一致收敛.求该级数的和函数 $F(x)$.
\end{example}
\exampleblank

\begin{solution}
    按题设可逐项求导,而求导只把双向级数的各项向左平移一位,所以
    \[
        F'(x)=F(x).
    \]
    故 $F(x)=Ce^x$.在 $x=0$ 处所有正重积分项为零,因此
    \[
        C=F(0)=\sum\limits_{n=0}^{\infty}f^{(n)}(0).
    \]
\end{solution}

%例题9.2.16
\begin{example}
    证明:存在 $\xi\in(0,1)$,使得
    \[
        \displaystyle\sum\limits_{n=1}^{\infty}\dfrac{\xi^{n-1}}{n+1}=1.
    \]
\end{example}
\exampleblank

\begin{solution}
    令
    \[
        F(x)=\sum\limits_{n=1}^{\infty}\dfrac{x^{n-1}}{n+1},\qquad0\leq x<1.
    \]
    该函数连续且严格递增,并且 $F(0)=1/2<1$,而
    $F(x)=[-\ln(1-x)-x]/x^2\to+\infty$($x\to1^-$).由介值定理,存在唯一 $\xi\in(0,1)$ 使 $F(\xi)=1$.
\end{solution}

%例题9.2.17
\begin{example}[\markstar]
    设 $f_n$ 均在 $\lbrack a,b\rbrack$ 上有原函数,且在 $\lbrack a,b\rbrack$ 上一致收敛于 $f(x)$,则 $f(x)$ 在 $\lbrack a,b\rbrack$ 上也有原函数.
\end{example}
\exampleblank

\begin{solution}
    取原函数 $F_n$ 并归一化为 $F_n(a)=0$.由 Lagrange 中值定理,
    \[
        \|F_n-F_m\|_\infty\leq(b-a)\|f_n-f_m\|_\infty.
    \]
    故 $F_n$ 一致收敛于某个 $F$.再由导数一致收敛定理,$F'=f$,所以 $f$ 有原函数.
\end{solution}

%例题9.2.18
\begin{example}[\markstar]
    设 $f_n$ 在 $\lbrack a,b\rbrack$ 上可积,且在 $\lbrack a,b\rbrack$ 上 $f_n(x)\rightrightarrows f(x)$,则 $f(x)$ 在 $\lbrack a,b\rbrack$ 上可积.
\end{example}
\exampleblank

\begin{solution}
    任取分割 $P$.由一致收敛,充分大时 $\|f_n-f\|_\infty<\varepsilon$,从而上下 Darboux 和满足
    \[
        U(f,P)-L(f,P)
        \leq U(f_n,P)-L(f_n,P)+2\varepsilon(b-a).
    \]
    先为可积函数 $f_n$ 选择使第一项任意小的分割,再令 $\varepsilon\to0$,即得 $f$ 可积.
\end{solution}


\newpage

\section{连续函数的逼近定理}

\begin{proposition}[\markstar\ Weierstrass 多项式逼近定理]
    有界闭区间 $\lbrack a,b\rbrack$ 上的连续函数 $f$ 一定可以用多项式一致逼近到任意程度,即对任意 $\varepsilon>0$,存在 $n=n(\varepsilon)$ 次多项式 $p_n$,使得
    \[
        |f(x)-p_n(x)|\leq\varepsilon.
    \]
\end{proposition}

\begin{proposition}[\markstar\ Weierstrass 三角多项式逼近定理]
    周期为 $2\pi$ 的连续函数 $f$ 一定可以用三角多项式一致逼近到任意程度,即对任意 $\varepsilon>0$,存在三角多项式
    \[
        S_n(x)=\dfrac{a_0}{2}+\sum\limits_{k=1}^{n}(a_k\cos kx+b_k\sin kx),
        \qquad n=n(\varepsilon),
    \]
    使得 $|f(x)-S_n(x)|<\varepsilon$.
\end{proposition}

%例题9.3.1
\begin{example}
    设 $f\in C\lbrack a,b\rbrack$,且对每个非负整数 $n$,满足
    \[
        \displaystyle\int_a^bx^nf(x)\,\mathrm{d}x=0.
    \]
    证明:$f\equiv0$.又若条件改为对大于某个正整数 $n_0$ 的所有 $n$ 成立,则也有相同结论.
\end{example}
\exampleblank

\begin{solution}
    由 Weierstrass 逼近定理,存在多项式 $p_m$ 一致逼近 $f$.题设给出
    $\displaystyle\int_a^bp_m(x)f(x)\,\mathrm{d}x=0$,故令 $m\to\infty$ 得
    \[
        \displaystyle\int_a^bf(x)^2\,\mathrm{d}x=0,
    \]
    从而 $f\equiv0$.

    若只对 $n>n_0$ 成立,则对函数 $x^{n_0+1}f(x)$ 的所有多项式矩仍为零,故
    $x^{n_0+1}f(x)\equiv0$;由连续性仍有 $f\equiv0$.
\end{solution}

%例题9.3.2
\begin{example}
    \begin{enumerate}
        \item 设 $f\in C\lbrack -1,1\rbrack$,且对每个非负整数 $n$ 满足 $\displaystyle\int_{-1}^{1}x^{2n}f(x)\,\mathrm{d}x=0$,则 $f$ 必为奇函数;
        \item 设 $f\in C\lbrack -1,1\rbrack$,且对每个非负整数 $n$ 满足 $\displaystyle\int_{-1}^{1}x^{2n+1}f(x)\,\mathrm{d}x=0$,则 $f$ 必为偶函数.
    \end{enumerate}
\end{example}
\exampleblank

\begin{solution}
    \leavevmode
    \begin{enumerate}
        \item 令 $f_e(x)=[f(x)+f(-x)]/2$ 为偶部.题设等价于
              $\displaystyle\int_{-1}^{1}x^{2n}f_e(x)\,\mathrm{d}x=0$.偶连续函数可由 $x^2$ 的多项式一致逼近,故与上一题同理得到 $f_e=0$,即 $f$ 为奇函数.
        \item 对奇部 $f_o(x)=[f(x)-f(-x)]/2$ 作同样论证;$x^{2n+1}f_o(x)$ 为偶函数,最终得到 $f_o=0$,故 $f$ 为偶函数.
    \end{enumerate}
\end{solution}

%例题9.3.3
\begin{example}
    若实系数多项式序列 $\{p_n(x)\}$ 在 $\mathbb{R}$ 上一致收敛于函数 $f(x)$,证明:$f(x)$ 必是多项式函数.
\end{example}
\exampleblank

\begin{solution}
    一致收敛使 $\{p_n\}$ 在 $\mathbb{R}$ 上一致 Cauchy.故充分大的 $m,n$ 满足
    $p_n-p_m$ 在整个实轴上有界.非恒定实多项式在实轴上无界,所以
    $p_n-p_m$ 必为常数.固定一个充分大的 $N$,便有 $p_n=p_N+c_n$,且
    $c_n$ 收敛.因此 $f=p_N+\displaystyle\lim c_n$,仍为多项式.
\end{solution}

%例题9.3.4
\begin{example}
    设 $f\in R\lbrack a,b\rbrack$,则对于每个 $\varepsilon>0$,存在两个多项式 $p(x)$ 和 $P(x)$,使得
    \begin{enumerate}
        \item $p(x)\leq f(x)\leq P(x)$, $x\in\lbrack a,b\rbrack$;
        \item $\displaystyle\int_a^b[P(x)-p(x)]\,\mathrm{d}x<\varepsilon$.
    \end{enumerate}
\end{example}
\exampleblank

\begin{solution}
    由 Riemann 可积性,可取分割使上下 Darboux 和之差小于 $\varepsilon/3$.把相应上下阶梯函数仅在各分点的充分小邻域内改成连续折线,可得连续函数 $u,v$ 满足
    \[
        u\leq f\leq v,\qquad \int_a^b(v-u)<\dfrac{2\varepsilon}{3}.
    \]
    再由 Weierstrass 定理分别用多项式一致逼近 $u$ 与 $v$,并把误差常数向下,向上平移,得到多项式 $p\leq u\leq f\leq v\leq P$,且
    $\displaystyle\int(P-p)<\varepsilon$.
\end{solution}

%例题9.3.5
\begin{example}
    设 $f\in C\lbrack 0,1\rbrack$,证明:存在每项为奇数项的多项式序列在 $\lbrack 0,1\rbrack$ 上一致收敛于 $f$ 的充分必要条件为 $f(0)=0$.
\end{example}
\exampleblank

\begin{solution}
    必要性显然,因为所有奇次幂多项式在零点均为零.

    充分性使用 M\"untz--Sz\'asz 定理.对指数列 $1,3,5,\ldots$,有
    \[
        \displaystyle\sum\limits_{k=0}^{\infty}\dfrac{1}{2k+1}=+\infty,
    \]
    故这些单项式的线性张成在所有满足 $f(0)=0$ 的 $C[0,1]$ 函数中稠密.因此存在只含奇次幂的多项式序列一致逼近 $f$.
\end{solution}

%例题9.3.6
\begin{example}
    设 $f\in R\lbrack a,b\rbrack$,且对每个非负整数 $n$ 满足 $\displaystyle\int_a^bx^nf(x)\,\mathrm{d}x=0$,证明:$f$ 在每个连续点上等于 $0$.
\end{example}
\exampleblank

\begin{solution}
    题设蕴含对任意多项式 $p$ 都有 $\displaystyle\int_a^bp(x)f(x)\,\mathrm{d}x=0$.若某连续点 $c$ 满足 $f(c)>0$,则在其某邻域内 $f$ 有正下界.选取在 $c$ 附近接近 $1$,在该邻域外一致趋于零的非负多项式峰函数 $p_N$,可使
    \[
        \displaystyle\int_a^bp_N(x)f(x)\,\mathrm{d}x>0,
    \]
    矛盾;$f(c)<0$ 同理.因此每个连续点处都有 $f(c)=0$.
\end{solution}

%例题9.3.7
\begin{example}
    设 $f\in C\lbrack 0,1\rbrack$,若
    \[
        \displaystyle\int_0^1f\left(x^{\frac{1}{2n+1}}\right)\,\mathrm{d}x=0,
    \]
    则 $f(x)\equiv0$.
\end{example}
\exampleblank

\begin{solution}
    对每个非负整数 $n$ 作代换 $x=t^{2n+1}$,题设成为
    \[
        (2n+1)\int_0^1t^{2n}f(t)\,\mathrm{d}t=0.
    \]
    而 $1,t^2,t^4,\ldots$ 的线性组合在 $C[0,1]$ 中稠密,因为任意
    $h(t)$ 可写成连续函数 $H(t^2)$.故可用这些多项式一致逼近 $f$,从而
    $\displaystyle\int_0^1f^2=0$,即 $f\equiv0$.
\end{solution}

%例题9.3.8
\begin{example}
    \begin{enumerate}
        \item 设 $f$ 在 $\lbrack a,b\rbrack$ 上有界且有原函数,$g\in C\lbrack a,b\rbrack$,则 $f(x)g(x)$ 在 $\lbrack a,b\rbrack$ 上有原函数;
        \item 设 $f(x)$ 在 $\lbrack a,b\rbrack$ 上有界且有原函数,$g:\lbrack 0,1\rbrack\to\lbrack a,b\rbrack$ 有连续导数且 $g'(x)>0$,则 $f[g(x)]$ 在 $\lbrack 0,1\rbrack$ 上有原函数.
    \end{enumerate}
\end{example}
\exampleblank

\begin{solution}
    \leavevmode
    \begin{enumerate}
        \item 设 $F'=f$.用多项式 $q_m$ 一致逼近 $g$.对每个多项式 $q$,
              \[
                  fq=(Fq)'-Fq'
              \]
              有原函数;又 $fq_m\rightrightarrows fg$,由前述一致极限定理,$fg$ 也有原函数.
        \item $g'>0$,故 $1/g'$ 连续.函数
              \[
                  f(g(x))g'(x)=\dfrac{\mathrm{d}}{\mathrm{d}x}F(g(x))
              \]
              有原函数.再对它与连续函数 $1/g'$ 应用 (1),即得 $f(g(x))$ 有原函数.
    \end{enumerate}
\end{solution}


\newpage

\section{幂级数}


\subsection{收敛半径与收敛域}

\begin{theorem}[\markstar]
    对于幂级数 $\displaystyle\sum\limits_{n=0}^{\infty}a_nx^n$:
    \begin{enumerate}
        \item 若该幂级数在 $x=x_1\neq0$ 处收敛,则对 $|x|<|x_1|$,幂级数绝对收敛;
        \item 若该幂级数在 $x=x_2$ 处发散,则对 $|x|>|x_2|$,幂级数发散.
    \end{enumerate}
\end{theorem}

\begin{corollary}[\markstar]
    当幂级数 $\displaystyle\sum\limits_{n=0}^{\infty}a_nx^n$ 具有收敛点 $x_1\neq0$,发散点 $x_2$ 时,必存在正数 $R$,使得 $|x|<R$ 时级数收敛,$|x|>R$ 时级数发散.
\end{corollary}

\begin{theorem}[\markstar]
    若
    \[
        \displaystyle\lim\limits_{n\to\infty}\left|\dfrac{a_{n+1}}{a_n}\right|=\rho,
    \]
    则收敛半径 $R=\dfrac{1}{\rho}$.进一步,
    \[
        R=\dfrac{1}{\displaystyle\lim\limits_{n\to\infty}\sqrt[n]{|a_n|}}
    \]
    在相应极限存在时成立.
\end{theorem}

\begin{proposition}[\markstar]
    \leavevmode
    \begin{enumerate}
        \item 设 $\{a_n\}$ 是有界列,则 $\displaystyle\sum\limits_{n=0}^{\infty}a_nx^n$ 的收敛半径 $R\geq1$.设该幂级数收敛半径 $R>0$,则对于 $|A|>R$,数列 $\{a_nA^n\}$ 无界;
        \item
              \[
                  R=\sup\{r\geq0\mid \{|a_n|r^n\}\text{ 是有界列}\}.
              \]
    \end{enumerate}
\end{proposition}

\begin{proposition}[\markstar]
    设 $\displaystyle\sum\limits_{n=1}^{\infty}a_nx^n$ 与 $\displaystyle\sum\limits_{n=1}^{\infty}b_nx^n$ 的收敛半径分别为正数 $R_1,R_2$,则
    \begin{enumerate}
        \item $\displaystyle\sum\limits_{n=0}^{\infty}a_nb_nx^n$ 的收敛半径 $R\geq R_1R_2$,而 $\displaystyle\sum\limits_{n=0}^{\infty}\dfrac{b_n}{a_n}x^n$ 的收敛半径 $R\leq\dfrac{R_1}{R_2}$;
        \item $R_1\neq R_2$ 时,$\displaystyle\sum\limits_{n=0}^{\infty}(a_n+b_n)x^n$ 的收敛半径 $R=\min\{R_1,R_2\}$;
        \item 两幂级数的 Cauchy 乘积的收敛半径 $R\geq\min\{R_1,R_2\}$.
    \end{enumerate}
\end{proposition}

\begin{proposition}[\markstar]
    设 $\displaystyle\sum\limits_{n=0}^{\infty}a_nx^n$ 的收敛半径为 $R$,则 $\displaystyle\sum\limits_{n=0}^{\infty}a_n^kx^n$ 的收敛半径为 $R^k$,而 $\displaystyle\sum\limits_{n=0}^{\infty}a_nx^{kn}$ 的收敛半径为 $\sqrt[k]{R}$.
\end{proposition}

%例题9.4.1
\begin{example}
    试求下列幂级数的收敛半径 $R$:
    \begin{enumerate}
        \item $\displaystyle\sum\limits_{n=1}^{\infty}\dfrac{(n!)^2}{(2n)!}x^n$;
        \item $\displaystyle\sum\limits_{n=1}^{\infty}\ln\dfrac{3n-2}{3n+2}x^n$;
        \item $\displaystyle\sum\limits_{n=1}^{\infty}\ln\left(\cos\dfrac{1}{3n}\right)x^n$;
        \item $\displaystyle\sum\limits_{n=1}^{\infty}\dfrac{1}{n!}\left(\dfrac{nx}{e}\right)^n$;
        \item $\displaystyle\sum\limits_{n=1}^{\infty}\dfrac{n^n}{n!}e^{-n\alpha}x^n$;
        \item $\displaystyle\sum\limits_{n=1}^{\infty}\dfrac{x^n}{a^n+b^n}$.
    \end{enumerate}
\end{example}
\exampleblank

\begin{solution}
    由 Cauchy--Hadamard 公式或比值判别法逐项计算,得到
    \[
        R_1=4,\quad R_2=1,\quad R_3=1,\quad R_4=1,
        \quad R_5=e^{\alpha-1},\quad R_6=\max\{|a|,|b|\}.
    \]
    其中第 (1) 问用 $(2n)!/(n!)^2\sim4^n/\sqrt{\pi n}$;第 (2)(3) 问分别用
    $\ln\frac{3n-2}{3n+2}\sim-\frac{4}{3n}$,$\ln(\cos\frac{1}{3n})\sim-\frac{1}{18n^2}$;第 (4)(5) 问用 Stirling 公式.第 (6) 问需假定 $a,b$ 不同时为零.
\end{solution}

%例题9.4.2
\begin{example}
    试求下列幂级数的收敛半径 $R$:
    \begin{enumerate}
        \item $\displaystyle\sum\limits_{n=0}^{\infty}5^n x^{3n}$;
        \item $\displaystyle\sum\limits_{n=1}^{\infty}n^n x^{n^2}$;
        \item $\displaystyle\sum\limits_{n=0}^{\infty}n!x^{n!}$.
    \end{enumerate}
\end{example}
\exampleblank

\begin{solution}
    把级数补成普通幂级数并使用 Cauchy--Hadamard 公式.非零系数分别位于 $3n,n^2,n!$ 次项,故
    \[
        R_1=5^{-1/3},\qquad R_2=1,qquad R_3=1.
    \]
    后两式分别利用 $(n^n)^{1/n^2}=n^{1/n}\to1$ 与 $(n!)^{1/n!}\to1$.
\end{solution}

%例题9.4.3
\begin{example}
    若存在 $\alpha>0,l>0$,使得
    \[
        \displaystyle\lim\limits_{n\to\infty}|a_nn^\alpha|=l,
    \]
    则 $\displaystyle\sum\limits_{n=0}^{\infty}a_nx^n$ 的收敛半径 $R=1$.
\end{example}
\exampleblank

\begin{solution}
    由 $|a_n|\sim l n^{-\alpha}$,有
    \[
        \displaystyle\lim\limits_{n\to\infty}\sqrt[n]{|a_n|}
        =\lim\limits_{n\to\infty}l^{1/n}n^{-\alpha/n}=1.
    \]
    故 Cauchy--Hadamard 公式给出 $R=1$.
\end{solution}

%例题9.4.4
\begin{example}
    设 $a_n>0$, $S_n=\displaystyle\sum\limits_{k=0}^{n}a_k$.若
    \[
        \displaystyle\lim\limits_{n\to\infty}S_n=+\infty,
        \qquad
        \lim\limits_{n\to\infty}\dfrac{a_n}{S_n}=0,
    \]
    则 $\displaystyle\sum\limits_{n=0}^{\infty}a_nx^n$ 的收敛半径 $R=1$.
\end{example}
\exampleblank

\begin{solution}
    由 $a_n/S_n\to0$ 得 $S_{n-1}/S_n\to1$.又
    \[
        \ln S_n-\ln S_{n-1}=\ln\dfrac{S_n}{S_{n-1}}\to0,
    \]
    由 Stolz 定理,$\ln S_n/n\to0$,故 $\sqrt[n]{S_n}\to1$.于是
    $\displaystyle\limsup\sqrt[n]{a_n}\leq1$,从而 $R\geq1$;而在 $x=1$ 时
    $\displaystyle\sum a_n=+\infty$,故 $R\leq1$.所以 $R=1$.
\end{solution}

%例题9.4.5
\begin{example}
    设 $a_n>0$, $\displaystyle\sum\limits_{n=0}^{\infty}a_nx^n$ 的收敛半径是 $1$,则
    \[
        I=\sum\limits_{n=0}^{\infty}S_nx^n,
        \qquad
        S_n=\sum\limits_{k=0}^{n}a_k,
    \]
    的收敛半径也是 $1$.
\end{example}
\exampleblank

\begin{solution}
    当 $|x|<1$ 时可交换求和次序:
    \[
        \displaystyle\sum\limits_{n=0}^{\infty}S_nx^n
        =\sum\limits_{k=0}^{\infty}a_k\sum\limits_{n=k}^{\infty}x^n
        =\dfrac{1}{1-x}\sum\limits_{k=0}^{\infty}a_kx^k.
    \]
    右端在 $|x|<1$ 收敛,故新级数半径至少为 $1$;因 $S_n\geq a_n>0$ 且原级数半径为 $1$,有
    $\displaystyle\limsup\sqrt[n]{S_n}\geq1$,故半径不超过 $1$.于是仍为 $1$.
\end{solution}

%例题9.4.6
\begin{example}
    证明:
    \[
        \displaystyle\sum\limits_{\substack{k+j=n\\0\leq k,j\leq n}}C_{2k}^{k}C_{2j}^{j}=4^n,
        \qquad n=0,1,2,\ldots.
    \]
\end{example}
\exampleblank

\begin{solution}
    利用母函数
    \[
        \displaystyle\sum\limits_{k=0}^{\infty}C_{2k}^{k}x^k=(1-4x)^{-1/2}.
    \]
    两边平方并比较 $x^n$ 的系数:
    \[
        \displaystyle\sum\limits_{k+j=n}C_{2k}^{k}C_{2j}^{j}
            =[x^n](1-4x)^{-1}=4^n.
    \]
\end{solution}

%例题9.4.7
\begin{example}
    试求下列幂级数的收敛域:
    \begin{enumerate}
        \item $\displaystyle\sum\limits_{n=1}^{\infty}\dfrac{x^n}{1+\dfrac{1}{2}+\cdots+\dfrac{1}{n}}$;
        \item $\displaystyle\sum\limits_{n=0}^{\infty}C_{\alpha}^{n}x^n$.
    \end{enumerate}
\end{example}
\exampleblank

\begin{solution}
    第 (1) 问中调和数 $H_n\sim\ln n$,故半径为 $1$;$x=1$ 时
    $\displaystyle\sum1/H_n$ 发散,$x=-1$ 时由 Leibniz 判别法收敛.因此收敛域为 $\lbrack-1,1)$.

    第 (2) 问按 $C_{\alpha}^{n}$ 表示广义二项式系数理解.若 $\alpha$ 为非负整数,级数退化为多项式,处处收敛;否则半径为 $1$.由
    \[
        C_{\alpha}^{n}\sim\dfrac{(-1)^n}{\Gamma(-\alpha)}n^{-\alpha-1}
    \]
    可再按 $\alpha$ 的取值判别端点:$x=1$ 在 $\alpha>-1$ 时收敛,$x=-1$ 在 $\alpha>0$ 时收敛;临界值需直接按有限项或调和型级数处理.
\end{solution}

%例题9.4.8
\begin{example}
    试求下列幂级数的收敛区域:
    \begin{enumerate}
        \item $\displaystyle\sum\limits_{n=1}^{\infty}(-1)^{n-1}\dfrac{2n+3}{3n^2+4}x^{2n+1}$;
        \item $\displaystyle\sum\limits_{n=1}^{\infty}(\sqrt[n]{a}-1)x^n$, $a>0$;
        \item $\displaystyle\sum\limits_{n=0}^{\infty}\sqrt{2-2\cos\dfrac{\pi}{2n}}\,x^n$;
        \item $\displaystyle\sum\limits_{n=1}^{\infty}\dfrac{x^n}{n^p\ln n}$, $p>0$;
        \item $\displaystyle\sum\limits_{n=1}^{\infty}\dfrac{[2+(-1)^n]^n}{n}x^n$;
        \item $\displaystyle\sum\limits_{n=1}^{\infty}\dfrac{(1+\dfrac{1}{n})(-1)^n}{n^2}x^n$.
    \end{enumerate}
\end{example}
\exampleblank

\begin{solution}
    各题半径及端点判别如下:
    \begin{enumerate}
        \item $|x|\leq1$,两端均绝对收敛;
        \item $|x|<1$,且 $x=-1$ 条件收敛,$x=1$ 发散,故 $\lbrack-1,1)$;
        \item 题面 $n=0$ 时无定义,应从 $n=1$ 起.此时为 $\lbrack-1,1)$;
        \item 题面 $n=1$ 时 $\ln n=0$,应从 $n=2$ 起.$x=-1$ 总收敛,$x=1$ 当且仅当 $p>1$ 收敛;
        \item $|x|<1/3$,在 $x=\pm1/3$ 均因偶数项形成调和型子级数而发散;
        \item $\lbrack-1,1\rbrack$,两端均绝对收敛.
    \end{enumerate}
\end{solution}

%例题9.4.9
\begin{example}
    试用变量替换,并借助幂级数法求下列函数项级数的收敛域:
    \begin{enumerate}
        \item $\displaystyle\sum\limits_{n=1}^{\infty}\dfrac{\sqrt[3]{2n+1}-\sqrt[3]{2n-1}}{\sqrt n}(x+3)^n$;
        \item $\displaystyle\sum\limits_{n=1}^{\infty}\dfrac{n+2}{2^n}e^{-nx}$;
        \item $\displaystyle\sum\limits_{n=1}^{\infty}\dfrac{1}{3+2n}\left(\dfrac{1-x}{1+x}\right)^n$;
        \item $\displaystyle\sum\limits_{n=1}^{\infty}\sqrt n(\tan x)^n$;
        \item $\displaystyle\sum\limits_{n=1}^{\infty}\dfrac{n4^n}{3^n}\dfrac{x^n}{(1-x)^n}$;
        \item $\displaystyle\sum\limits_{n=1}^{\infty}\dfrac{2^{-n}}{\sqrt n}x^n$.
    \end{enumerate}
\end{example}
\exampleblank

\begin{solution}
    分别令括号中的量为新变量并判别相应幂级数,得到
    \begin{enumerate}
        \item 系数为 $O(n^{-7/6})$,故 $|x+3|\leq1$,即 $\lbrack-4,-2\rbrack$;
        \item $|e^{-x}/2|<1$,故 $x>-\ln2$;
        \item $|(1-x)/(1+x)|<1$,故 $x>0$;
        \item $|\tan x|<1$,即各区间 $k\pi-\pi/4<x<k\pi+\pi/4$;
        \item $|4x/[3(1-x)]|<1$,故 $-3<x<3/7$;
        \item 半径为 $2$,$x=-2$ 条件收敛而 $x=2$ 发散,故 $\lbrack-2,2)$.
    \end{enumerate}
\end{solution}

%例题9.4.10
\begin{example}
    设 $a_0=0,a_1=1,a_{n+1}=a_{n-1}+a_n$,则 $\displaystyle\sum\limits_{n=0}^{\infty}a_nx^n$ 在 $|x|<\dfrac{1}{2}$ 时收敛.
\end{example}
\exampleblank

\begin{solution}
    由递推式归纳可得 $0<a_n\leq2^{n-1}$ $(n\geq1)$.因此当 $|x|<1/2$ 时
    \[
        \displaystyle\sum\limits_{n=0}^{\infty}|a_nx^n|
        \leq\sum\limits_{n=1}^{\infty}2^{n-1}|x|^n<\infty.
    \]
    事实上母函数为 $x/(1-x-x^2)$,精确收敛半径为 $(\sqrt5-1)/2$.
\end{solution}

%例题9.4.11
\begin{example}
    若 $\displaystyle\sum\limits_{n=0}^{\infty}a_n^2x^n$ 的收敛域为 $\lbrack -1,1\rbrack$,则 $\displaystyle\sum\limits_{n=0}^{\infty}\dfrac{a_n}{n}x^n$ 的收敛域为 $\lbrack -1,1\rbrack$.
\end{example}
\exampleblank

\begin{solution}
    题面第二个级数的 $n=0$ 项无定义,应从 $n=1$ 起.由第一级数半径为 $1$,
    \[
        \displaystyle\limsup\limits_{n\to\infty}|a_n|^{2/n}=1,
    \]
    故 $\displaystyle\limsup|a_n/n|^{1/n}=1$,第二级数半径也为 $1$.又 $\displaystyle\sum a_n^2$ 收敛,由 Cauchy--Schwarz 不等式
    \[
        \displaystyle\sum\limits_{n=1}^{\infty}\dfrac{|a_n|}{n}
        \leq\left(\sum a_n^2\right)^{1/2}
        \left(\sum\dfrac{1}{n^2}\right)^{1/2}<\infty.
    \]
    故在 $x=\pm1$ 均绝对收敛,收敛域为 $\lbrack-1,1\rbrack$.
\end{solution}

%例题9.4.12
\begin{example}
    若 $\displaystyle\sum\limits_{n=1}^{\infty}a_n$ 在 Cesaro 意义下可求和,则 $\displaystyle\sum\limits_{n=1}^{\infty}a_nx^n$ 在 $(-1,1)$ 收敛.
\end{example}
\exampleblank

\begin{solution}
    记 $s_n=\displaystyle\sum_{k=1}^na_k$,Cesaro 可求和意味着
    $\sigma_n=(s_1+\cdots+s_n)/n$ 有有限极限.于是
    \[
        \dfrac{s_n}{n}=\dfrac{n\sigma_n-(n-1)\sigma_{n-1}}{n}\longrightarrow0,
    \]
    从而 $a_n=s_n-s_{n-1}=o(n)$.对任意 $|x|<1$,级数
    $\displaystyle\sum |a_nx^n|$ 被 $\displaystyle\sum Cn|x|^n$ 控制,故绝对收敛.
\end{solution}



\subsection{幂级数求和}

\begin{theorem}[Tauber 定理 $*$]
    设
    \[
        f(x)=\sum\limits_{n=0}^{\infty}a_nx^n,
        \qquad -1<x<1,
    \]
    且 $na_n\to0$.若 $\displaystyle\lim\limits_{x\to1^-}f(x)=S$,则
    \[
        \displaystyle\sum\limits_{n=0}^{\infty}a_n=S.
    \]
\end{theorem}

\begin{remark}
    Abel 第二定理的逆形式一般并不成立.例如
    \[
        f(x)=\dfrac{1}{1+x}=\sum\limits_{n=0}^{\infty}(-1)^nx^n,
        \qquad -1<x<1,
    \]
    有 $f(x)\to\dfrac{1}{2}$ $(x\to1^-)$,但 $\displaystyle\sum\limits_{n=0}^{\infty}(-1)^n$ 发散.Tauber 定理表明,对系数增加 $na_n\to0$ 的限制后,逆命题成立.
\end{remark}

幂级数的求和方法主要有两种:逐项积分与求导,Abel 第二定理.

%例题9.4.13
\begin{example}
    设 $a_n>0$.若
    \[
        \displaystyle\lim\limits_{x\to1^-}\sum\limits_{n=0}^{\infty}a_nx^n=S,
    \]
    则 $\displaystyle\sum\limits_{n=0}^{\infty}a_n=S$.
\end{example}
\exampleblank

\begin{solution}
    因 $a_n>0$,部分和 $S_N=\displaystyle\sum_{n=0}^Na_n$ 单调增加.对任意固定 $N$,
    \[
        \displaystyle\sum\limits_{n=0}^Na_nx^n\leq\sum\limits_{n=0}^{\infty}a_nx^n.
    \]
    令 $x\to1^-$ 得 $S_N\leq S$.再令 $N\to\infty$,得到
    $\displaystyle\sum a_n\leq S$.反向不等式由 $x^n\leq1$ 得到,故 $\displaystyle\sum a_n=S$.
\end{solution}

%例题9.4.14
\begin{example}[Tauber 定理的推广 $*$]
    设 $f(x)=\displaystyle\sum\limits_{n=0}^{\infty}a_nx^n$ 的收敛半径为 $1$.若
    \[
        \displaystyle\lim\limits_{n\to\infty}\dfrac{1}{n}\sum\limits_{k=1}^{n}ka_k=0,
        \qquad
        \lim\limits_{x\to1^-}f(x)=l,
    \]
    则 $\displaystyle\sum\limits_{n=0}^{\infty}a_n=l$.
\end{example}
\exampleblank

\begin{solution}
    记 $s_n=\displaystyle\sum_{k=0}^na_k$.分部求和可把 $s_n-l$ 写成 Abel 平均误差与
    $n^{-1}\displaystyle\sum_{k=1}^nka_k$ 的线性组合.具体取 $x=1-1/n$,有
    \[
        s_n-f(x)=\sum\limits_{k=0}^na_k(1-x^k)-\sum\limits_{k>n}a_kx^k.
    \]
    对两项作 Abel 变换,并用
    $\frac{1}{n}\displaystyle\sum_{k=1}^nka_k\to0$,可得右端趋于零.又 $f(1-1/n)\to l$,故 $s_n\to l$.这正是题设 Tauber 条件下的结论.
\end{solution}

\begin{remark}
    若将 $\displaystyle\dfrac{1}{n}\sum\limits_{k=1}^{n}ka_k\to0$ 换成 $\displaystyle\sum\limits_{n=1}^{\infty}na_n^2$ 收敛,结论依旧成立.
\end{remark}

%例题9.4.15
\begin{example}
    求下列幂级数的和 $S(x)$:
    \begin{enumerate}
        \item $\displaystyle\sum\limits_{n=1}^{\infty}nx^n$, $|x|<1$;
        \item $\displaystyle\sum\limits_{n=1}^{\infty}\dfrac{x^n}{n}$, $-1\leq x<1$;
        \item $\displaystyle\sum\limits_{n=1}^{\infty}\dfrac{x^n}{n(n+1)(n+2)}$, $|x|\leq1$;
        \item $\displaystyle\sum\limits_{n=1}^{\infty}\dfrac{x^{2n-1}}{2n-1}$, $|x|<1$.
    \end{enumerate}
\end{example}
\exampleblank

\begin{solution}
    利用几何级数逐项微分或积分,得到
    \begin{enumerate}
        \item $\displaystyle S(x)=\dfrac{x}{(1-x)^2}$;
        \item $S(x)=-\ln(1-x)$,并在 $x=-1$ 取值 $-\ln2$;
        \item 令 $A=-\ln(1-x)$,则
              \[
                  S(x)=\dfrac{1}{2}\left[A-\dfrac{2(A-x)}x+dfrac{A-x-x^2/2}{x^2}\right],
              \]
              在 $x=0$ 处按连续延拓取值;
        \item $\displaystyle S(x)=\operatorname{arth}x
                  =\dfrac{1}{2}\ln\dfrac{1+x}{1-x}$.
    \end{enumerate}
\end{solution}

%例题9.4.16
\begin{example}
    求下列幂级数的和 $S(x)$:
    \begin{enumerate}
        \item $1+2x-4x^3-5x^4+7x^6+8x^7-\cdots$, $|x|<1$;
        \item $\displaystyle\sum\limits_{n=1}^{\infty}n^3x^n$, $|x|<1$;
        \item $\displaystyle\sum\limits_{n=1}^{\infty}\dfrac{(-1)^{n-1}x^{2n}}{n(2n-1)}$, $|x|<1$;
        \item 设 $\displaystyle S(x)=\sum\limits_{n=1}^{\infty}\dfrac{x^n}{n^2}$,求 $S(x)+S(1-x)$, $0<x<1$.
    \end{enumerate}
\end{example}
\exampleblank

\begin{solution}
    令 $y=-x^3$.各题求和为
    \begin{enumerate}
        \item $\displaystyle \dfrac{3y}{(1-y)^2}+\dfrac{1}{1-y}
                  +x\left[\dfrac{3y}{(1-y)^2}+\dfrac{2}{1-y}\right]$;
        \item $\displaystyle\dfrac{x(1+4x+x^2)}{(1-x)^4}$;
        \item $2x\arctan x-\ln(1+x^2)$;
        \item $\displaystyle S(x)+S(1-x)=\dfrac{\pi^2}{6}-\ln x\ln(1-x)$.
    \end{enumerate}
    最后一式可对左端求导并在 $x=1/2$ 处定常数.
\end{solution}

%例题9.4.17
\begin{example}
    试用幂级数法求下列常数项级数的和 $S$:
    \begin{enumerate}
        \item $\displaystyle S=\sum\limits_{n=2}^{\infty}\dfrac{(-1)^n}{n^2+n-2}$;
        \item $\displaystyle S=\sum\limits_{n=1}^{\infty}\dfrac{1}{n(2n+1)2^n}$;
        \item $\displaystyle S=\sum\limits_{n=0}^{\infty}\dfrac{(-1)^n}{3n+1}$;
        \item $\displaystyle S=1-\dfrac{1}{5}+\dfrac{1}{7}-\dfrac{1}{11}+\cdots+\dfrac{1}{6n-5}-\dfrac{1}{6n-1}+\cdots$.
    \end{enumerate}
\end{example}
\exampleblank

\begin{solution}
    由几何级数积分并化简,得到
    \begin{enumerate}
        \item $\displaystyle -\dfrac{5}{18}+\dfrac{2}{3}\ln2$;
        \item $\displaystyle 2+\ln2-2\sqrt2\operatorname{arth}\dfrac{1}{\sqrt2}$;
        \item $\displaystyle \int_0^1\dfrac{\mathrm{d}x}{1+x^3}
                  =\dfrac{1}{3}\ln2+\dfrac{\pi}{3\sqrt3}$;
        \item $\displaystyle\dfrac{\pi}{2\sqrt3}$.
    \end{enumerate}
\end{solution}

%例题9.4.18
\begin{example}
    试用公式
    \[
        \dfrac{1}{k(k+1)\cdots(k+n)}
        =\dfrac{1}{n!}\int_0^1(1-x)^nx^{k-1}\,\mathrm{d}x
    \]
    求下列级数的和:
    \begin{enumerate}
        \item $\displaystyle\dfrac{1}{1\cdot2\cdot3}+\dfrac{1}{3\cdot4\cdot5}+\dfrac{1}{5\cdot6\cdot7}+\cdots$;
        \item $\displaystyle\dfrac{1}{1\cdot2\cdot3\cdot4}+\dfrac{1}{4\cdot5\cdot6\cdot7}+\dfrac{1}{7\cdot8\cdot9\cdot10}+\cdots$.
    \end{enumerate}
\end{example}
\exampleblank

\begin{solution}
    把题给积分公式分别用于 $k=2m+1$ 与 $k=3m+1$,并交换求和与积分:
    \begin{enumerate}
        \item $\displaystyle\dfrac{1}{2}\int_0^1\dfrac{(1-x)^2}{1-x^2}\,\mathrm{d}x
                  =\ln2-\dfrac{1}{2}$;
        \item $\displaystyle\dfrac{1}{6}\int_0^1\dfrac{(1-x)^3}{1-x^3}\,\mathrm{d}x
                  =\dfrac{1}{6}-\dfrac{1}{4}\ln3+\dfrac{\sqrt3\pi}{36}$.
    \end{enumerate}
\end{solution}

%例题9.4.19
\begin{example}
    计算极限
    \[
        \displaystyle\lim\limits_{\substack{n\to\infty\\m\to\infty}}
        \sum\limits_{i=1}^{m}\sum\limits_{j=1}^{n}\dfrac{(-1)^{i+j}}{i+j}.
    \]
\end{example}
\exampleblank

\begin{solution}
    利用 $1/(i+j)=\displaystyle\int_0^1x^{i+j-1}\,\mathrm{d}x$,有限双和等于
    \[
        \displaystyle\int_0^1\dfrac{1}{x}
        \left(\sum\limits_{i=1}^m(-x)^i\right)
        \left(\sum\limits_{j=1}^n(-x)^j\right)\mathrm{d}x.
    \]
    令 $m,n\to\infty$,由控制收敛定理得到
    \[
        \displaystyle\int_0^1\dfrac{x}{(1+x)^2}\,\mathrm{d}x=\ln2-\dfrac{1}{2}.
    \]
\end{solution}

%例题9.4.20
\begin{example}
    设 $\displaystyle\lim\limits_{n\to\infty}a_n=l$,则 $\displaystyle S(x)=\sum\limits_{n=1}^{\infty}a_nx^n$ 在 $(-1,1)$ 上收敛,且
    \[
        \displaystyle\lim\limits_{x\to1^-}(1-x)S(x)=l.
    \]
\end{example}
\exampleblank

\begin{solution}
    写成 $a_n=l+(a_n-l)$.常数部分满足
    $(1-x)\displaystyle\sum lx^n=lx\to l$.对余项,给定 $\varepsilon>0$,取 $N$ 使
    $|a_n-l|<\varepsilon$ $(n>N)$,则
    \[
        (1-x)\left|\sum\limits_{n>N}(a_n-l)x^n\right|\leq\varepsilon.
    \]
    有限个首项乘 $1-x$ 后趋于零,故结论成立.
\end{solution}

%例题9.4.21
\begin{example}
    设 $a_n>0,b_n>0$,且
    \[
        S_1(x)=\sum\limits_{n=0}^{\infty}a_nx^n,
        \qquad
        S_2(x)=\sum\limits_{n=0}^{\infty}b_nx^n
    \]
    的收敛半径均为 $1$.如果
    \[
        \displaystyle\lim\limits_{x\to1^-}S_1(x)=\lim\limits_{x\to1^-}S_2(x)=+\infty,
        \qquad
        \lim\limits_{n\to\infty}\dfrac{a_n}{b_n}=l,
    \]
    则
    \[
        \displaystyle\lim\limits_{x\to1^-}\dfrac{S_1(x)}{S_2(x)}=l.
    \]
\end{example}
\exampleblank

\begin{solution}
    给定 $\varepsilon>0$,取 $N$ 使 $|a_n/b_n-l|<\varepsilon$ $(n>N)$.于是尾项满足
    \[
        (l-\varepsilon)\sum_{n>N}b_nx^n
        \leq\sum_{n>N}a_nx^n
        \leq(l+\varepsilon)\sum_{n>N}b_nx^n.
    \]
    因 $S_2(x)\to+\infty$,有限个首项除以 $S_2(x)$ 后趋于零.令 $x\to1^-$ 再令 $\varepsilon\to0$ 即得比值极限为 $l$.
\end{solution}

%例题9.4.22
\begin{example}[\markstar]
    设幂级数 $\displaystyle\sum\limits_{n=0}^{\infty}a_nx^n$ 的系数满足
    \[
        a_n+Aa_{n-1}+Ba_{n-2}=0,
    \]
    且此幂级数在点 $x=x_0$ 处收敛.证明:
    \[
        S=\sum\limits_{n=0}^{\infty}a_nx_0^n
        =\dfrac{a_0+(a_1+Aa_0)x_0}{1+Ax_0+Bx_0^2}.
    \]
\end{example}
\exampleblank

\begin{solution}
    把递推式乘以 $x_0^n$ 并从 $n=2$ 起求和.由于原幂级数在 $x_0$ 收敛,可作指标平移,得到
    \[
        (1+Ax_0+Bx_0^2)S=a_0+(a_1+Aa_0)x_0.
    \]
    若分母为零,则左端恒等式同时迫使分子为零,需另行按收敛条件讨论;在题设公式有意义时相除即得结论.
\end{solution}

%例题9.4.23
\begin{example}[\markstar\ 母函数法]
    设 $a_0=1,a_1=0,a_2=-5$,
    \[
        a_n=4a_{n-1}-5a_{n-2}+2a_{n-3}.
    \]
    证明:
    \[
        a_n=-2^{n+2}+3n+5.
    \]
\end{example}
\exampleblank

\begin{solution}
    直接验证即可.令 $u_n=-2^{n+2}+3n+5$,则
    $u_0=1,u_1=0,u_2=-5$,并且
    \[
        u_n=4u_{n-1}-5u_{n-2}+2u_{n-3}.
    \]
    线性递推在给定前三项后解唯一,故 $a_n=u_n$.
\end{solution}



\subsection{幂级数展开----Taylor 级数}

\begin{theorem}[\marktwostar\ 常用初等函数的 Maclaurin 级数展式]
    \begin{align*}
        \sin x                  & =x-\dfrac{x^3}{3!}+\dfrac{x^5}{5!}-\cdots+(-1)^n\dfrac{x^{2n+1}}{(2n+1)!}+\cdots, \\
        \cos x                  & =1-\dfrac{x^2}{2!}+\dfrac{x^4}{4!}-\cdots+(-1)^n\dfrac{x^{2n}}{(2n)!}+\cdots,     \\
        \tan x                  & =x+\dfrac{x^3}{3}+\dfrac{2x^5}{15}+\dfrac{17x^7}{315}+\cdots,                     \\
        e^x                     & =1+x+\dfrac{x^2}{2!}+\cdots+\dfrac{x^n}{n!}+\cdots,                               \\
        \ln(1+x)                & =x-\dfrac{x^2}{2}+\dfrac{x^3}{3}-\cdots+(-1)^{n+1}\dfrac{x^n}{n}+\cdots,          \\
        \arcsin x               & =x+\sum\limits_{n=1}^{\infty}\dfrac{(2n-1)!!}{(2n)!!}\dfrac{x^{2n+1}}{2n+1},      \\
        \arctan x               & =x-\dfrac{x^3}{3}+\dfrac{x^5}{5}-\cdots+(-1)^n\dfrac{x^{2n+1}}{2n+1}+\cdots,      \\
        (1+x)^\alpha            & =\sum\limits_{n=0}^{\infty}C_{\alpha}^{n}x^n,                                     \\
        \dfrac{1}{\sqrt{1-x^2}} & =1+\sum\limits_{n=1}^{\infty}\dfrac{(2n-1)!!}{(2n)!!}x^{2n}.
    \end{align*}
    相应收敛区间与原书一致.
\end{theorem}

\begin{remark}
    若 $f(x)$ 在 $x=x_0$ 处可幂级数展开,则
    \[
        f(x)=\sum\limits_{n=0}^{\infty}a_n(x-x_0)^n
        =\sum\limits_{n=0}^{\infty}\dfrac{f^{(n)}(x_0)}{n!}(x-x_0)^n,
    \]
    即为 Taylor 级数.幂级数展开常见方法有三种:分解合成法,微分积分法,待定系数法.
\end{remark}

%例题9.4.24
\begin{example}
    求下列函数在 $x=0$ 处的 Taylor 级数展式:
    \begin{enumerate}
        \item $f(x)=\dfrac{3x+8}{(2x-3)(x^2+4)}$;
        \item $f(x)=\ln(4+3x-x^2)$;
        \item $f(x)=\sin^4x$.
    \end{enumerate}
\end{example}
\exampleblank

\begin{solution}
    分解并使用基本级数:
    \begin{enumerate}
        \item
              \[
                  \dfrac{3x+8}{(2x-3)(x^2+4)}
                  =-\dfrac{2}{3}\sum\limits_{n=0}^{\infty}\left(\dfrac{2x}{3}\right)^n
                  -\sum\limits_{n=0}^{\infty}\dfrac{(-1)^nx^{2n+1}}{4^{n+1}},quad |x|<\dfrac{3}{2}.
              \]
        \item
              \[
                  \ln(4+3x-x^2)=\ln4+\sum\limits_{n=1}^{\infty}
                  \dfrac{(-1)^{n-1}-4^{-n}}n x^n,quad |x|<1.
              \]
        \item $\displaystyle\sin^4x=(3-4\cos2x+\cos4x)/8$,再展开两个余弦即可.
    \end{enumerate}
\end{solution}

%例题9.4.25
\begin{example}
    求下列函数的 Maclaurin 展开式:
    \begin{enumerate}
        \item $f(x)=\ln(1+x+x^2+x^3)$;
        \item $f(x)=\ln^2(1+x)$;
        \item $f(x)=\dfrac{\ln(1+x)}{1+x}$.
    \end{enumerate}
\end{example}
\exampleblank

\begin{solution}
    由基本对数级数及 Cauchy 乘积,
    \begin{enumerate}
        \item $\displaystyle\ln(1+x+x^2+x^3)=
                  \sum\limits_{n=1}^{\infty}\dfrac{(-1)^{n-1}x^n}{n}
                  +\sum\limits_{n=1}^{\infty}\dfrac{(-1)^{n-1}x^{2n}}n$;
        \item $\displaystyle\ln^2(1+x)=
                  2\sum\limits_{n=2}^{\infty}\dfrac{(-1)^nH_{n-1}}n x^n$;
        \item $\displaystyle\dfrac{\ln(1+x)}{1+x}=
                  \sum\limits_{n=1}^{\infty}(-1)^{n-1}H_nx^n$.
    \end{enumerate}
    上述展开在 $|x|<1$ 成立.
\end{solution}

%例题9.4.26
\begin{example}
    利用待定系数法求下列函数的 Maclaurin 级数展开式:
    \begin{enumerate}
        \item $f(x)=\dfrac{x\sin\alpha}{1-2x\cos\alpha+x^2}$;
        \item $f(x)=\dfrac{\arcsin x}{\sqrt{1-x^2}}$;
        \item $f(x)=\dfrac{1}{\sqrt{1-2ax+x^2}}$.
    \end{enumerate}
\end{example}
\exampleblank

\begin{solution}
    比较系数可得
    \begin{enumerate}
        \item $\displaystyle\sum\limits_{n=1}^{\infty}\sin(n\alpha)x^n$;
        \item $\displaystyle\dfrac{1}{2}\sum\limits_{n=1}^{\infty}
                  \dfrac{4^n}{nC_{2n}^{n}}x^{2n-1}$;
        \item $\displaystyle\sum\limits_{n=0}^{\infty}P_n(a)x^n$,其中 $P_n$ 为 Legendre 多项式.
    \end{enumerate}
    前两式由几何级数和 $(\arcsin x)^2$ 的展开求导得到;第三式正是 Legendre 多项式的生成函数,均在距最近奇点以内成立.
\end{solution}

%例题9.4.27
\begin{example}
    利用微分积分法求下列函数的 Maclaurin 级数展开式:
    \begin{enumerate}
        \item $f(x)=\arctan\dfrac{x+3}{x-3}$;
        \item $f(x)=\ln\left(x+\sqrt{x^2+1}\right)$;
        \item $f(x)=\arccos(1-2x^2)$.
    \end{enumerate}
\end{example}
\exampleblank

\begin{solution}
    逐项微分后积分回去,得到
    \begin{enumerate}
        \item $\displaystyle-\dfrac{\pi}{4}-
                  \sum\limits_{n=0}^{\infty}\dfrac{(-1)^n}{3\cdot9^n(2n+1)}x^{2n+1}$,$|x|<3$;
        \item $\displaystyle\sum\limits_{n=0}^{\infty}
                  \dfrac{(-1)^nC_{2n}^{n}}{4^n(2n+1)}x^{2n+1}$,$|x|<1$;
        \item 题面函数等于 $2\arcsin|x|$,在 $x=0$ 不可微,故不存在双侧 Maclaurin 级数.若限于 $x\geq0$,则为
              $\displaystyle2\sum_{n=0}^{\infty}\frac{C_{2n}^{n}}{4^n(2n+1)}x^{2n+1}$.
    \end{enumerate}
\end{solution}



\subsection{幂级数展开式的应用}


\methodsection{(A) 求级数的和}

%例题9.4.28
\begin{example}
    试求下列级数的和:
    \begin{enumerate}
        \item $\displaystyle S(x)=\sum\limits_{n=0}^{\infty}\dfrac{\ln^n x}{n!}$, $x>0$;
        \item $\displaystyle S(x)=\sum\limits_{n=0}^{\infty}\dfrac{(-1)^n\ln^n x}{2^n n!}$;
        \item $\displaystyle S(x)=\sum\limits_{n=0}^{\infty}\dfrac{2^n(n+1)}{n!}$.
    \end{enumerate}
\end{example}
\exampleblank

\begin{solution}
    由指数级数直接得到
    \[
        S_1(x)=e^{\ln x}=x,qquad
        S_2(x)=e^{-\frac{1}{2}\ln x}=x^{-1/2},qquad
        S_3=\sum\limits_{n=0}^{\infty}\dfrac{(n+1)2^n}{n!}=3e^2.
    \]
\end{solution}

%例题9.4.29
\begin{example}
    设 $\dfrac{1}{1-x-x^2}$ 的 Maclaurin 级数为 $\displaystyle\sum\limits_{n=0}^{\infty}a_nx^n$,试证明:
    \[
        \displaystyle\sum\limits_{n=0}^{\infty}\dfrac{a_{n+1}}{a_na_{n+2}}=2.
    \]
\end{example}
\exampleblank

\begin{solution}
    系数满足 Fibonacci 递推 $a_{n+2}=a_{n+1}+a_n$,且 $a_0=a_1=1$.因此
    \[
        \dfrac{a_{n+1}}{a_na_{n+2}}
        =\dfrac{1}{a_n}-\dfrac{1}{a_{n+2}}.
    \]
    求和至 $N$ 后望远镜相消,再令 $N\to\infty$,余项趋于零,所得和为
    $1/a_0+1/a_1=2$.
\end{solution}



\methodsection{(B) 求积分值}

%例题9.4.30
\begin{example}
    求下列积分值:
    \begin{enumerate}
        \item $\displaystyle\int_0^1e^x\ln x\,\mathrm{d}x$;
        \item $\displaystyle\int_0^{+\infty}\dfrac{x}{e^{2\pi x}-1}\,\mathrm{d}x$;
        \item $\displaystyle\int_0^x\dfrac{\sin t}{t}\,\mathrm{d}t$;
        \item $\displaystyle\int_0^1\ln x\ln(1-x)\,\mathrm{d}x$.
    \end{enumerate}
\end{example}
\exampleblank

\begin{solution}
    逐项展开并积分:
    \begin{enumerate}
        \item $\displaystyle-\sum\limits_{n=0}^{\infty}\dfrac{1}{n!(n+1)^2}$;
        \item $\displaystyle\sum\limits_{n=1}^{\infty}\dfrac{1}{4\pi^2n^2}=\dfrac{1}{24}$;
        \item $\displaystyle\sum\limits_{n=0}^{\infty}\dfrac{(-1)^nx^{2n+1}}{(2n+1)(2n+1)!}$;
        \item $\displaystyle\sum\limits_{n=1}^{\infty}\dfrac{1}{n(n+1)^2}=2-\dfrac{\pi^2}{6}$.
    \end{enumerate}
\end{solution}

%例题9.4.31
\begin{example}
    证明下列等式:
    \begin{enumerate}
        \item $\displaystyle\int_0^1\dfrac{\ln(1+x)}{x}\,\mathrm{d}x=\dfrac{\pi^2}{12}$;
        \item $\displaystyle\int_0^{\frac{\pi}{2}}\dfrac{\mathrm{d}\theta}{\sqrt{1-k^2\sin^2\theta}}$, $0\leq k<1$.
    \end{enumerate}
\end{example}
\exampleblank

\begin{solution}
    第 (1) 问由 $\ln(1+x)=\displaystyle\sum_{n\geq1}(-1)^{n-1}x^n/n$ 逐项除以 $x$ 并积分,得到交错平方和 $\pi^2/12$.

    第 (2) 问为第一类完全椭圆积分.由二项式级数
    \[
        (1-u)^{-1/2}=\sum\limits_{n=0}^{\infty}\dfrac{C_{2n}^{n}}{4^n}u^n
    \]
    及 $\displaystyle\int_0^{\pi/2}\sin^{2n}\theta\,\mathrm{d}\theta
        =\frac{\pi}{2}\frac{C_{2n}^{n}}{4^n}$,得
    \[
        \displaystyle\int_0^{\pi/2}\dfrac{\mathrm{d}\theta}{\sqrt{1-k^2\sin^2\theta}}
        =\dfrac{\pi}{2}\sum\limits_{n=0}^{\infty}
        \left(\dfrac{C_{2n}^{n}}{4^n}\right)^2k^{2n}.
    \]
\end{solution}



\methodsection{(C) 求导数 \texorpdfstring{$f^{(n)}(0)$}{f(n)(0)}}

%例题9.4.32
\begin{example}
    设
    \[
        f(x)=\dfrac{1-x+x^2}{1+x+x^2},
    \]
    试求 $f^{(n)}(0)$.
\end{example}
\exampleblank

\begin{solution}
    长除法或递推展开得
    \[
        f(x)=1+2\sum\limits_{k=0}^{\infty}(x^{3k+2}-x^{3k+1}).
    \]
    因此
    \[
        f^{(0)}(0)=1,qquad
        f^{(n)}(0)=
        \begin{cases}
            -2n!, & n\equiv1\pmod3,         \\
            2n!,  & n\equiv2\pmod3,         \\
            0,    & n\equiv0\pmod3, n\geq3.
        \end{cases}
    \]
\end{solution}


\newpage

\section{Fourier 级数}


\subsection{Fourier 系数}

设 $f(x)$ 以 $2\pi$ 为周期,在区间 $\lbrack -\pi,\pi\rbrack$ 上可积,则
\[
    a_n=\dfrac{1}{\pi}\int_{-\pi}^{\pi}f(x)\cos nx\,\mathrm{d}x,
    \qquad
    b_n=\dfrac{1}{\pi}\int_{-\pi}^{\pi}f(x)\sin nx\,\mathrm{d}x.
\]
称为 $f(x)$ 的 Fourier 系数,
\[
    \dfrac{a_0}{2}+\sum\limits_{n=1}^{\infty}(a_n\cos nx+b_n\sin nx)
\]
称为 $f(x)$ 在 $\lbrack -\pi,\pi\rbrack$ 上的 Fourier 级数.

\begin{proposition}[\markstar]
    Fourier 系数具有如下性质:
    \begin{enumerate}
        \item 线性性:若 $f_1,f_2$ 的 Fourier 系数分别为 $a_n^{(1)},b_n^{(1)}$ 与 $a_n^{(2)},b_n^{(2)}$,则 $\alpha f_1\pm\beta f_2$ 的系数为
              \[
                  a_n=\alpha a_n^{(1)}\pm\beta a_n^{(2)},
                  \qquad
                  b_n=\alpha b_n^{(1)}\pm\beta b_n^{(2)}.
              \]
        \item 若 $f$ 在 $\lbrack -\pi,\pi\rbrack$ 上连续,分段光滑,则 $f'$ 的 Fourier 系数满足 $b_n'=-na_n$;当 $f(-\pi)=f(\pi)$ 时,$a_0'=0$, $a_n'=nb_n$.
        \item 若 $f$ 在 $\lbrack -\pi,\pi\rbrack$ 上分段连续,设 $a_n,b_n$ 为其 Fourier 系数,则
              \[
                  F(x)=\int_0^x\left(f(t)-\dfrac{a_0}{2}\right)\mathrm{d}t
              \]
              的 Fourier 系数为
              \[
                  A_n=-\dfrac{b_n}{n},
                  \qquad
                  B_n=\dfrac{a_n}{n},
                  \qquad
                  A_0=2\sum\limits_{n=1}^{\infty}\dfrac{b_n}{n}.
              \]
    \end{enumerate}
\end{proposition}

\begin{proposition}[\markstar]
    设 $f(x)$ 以 $2\pi$ 为周期,$\lbrack -\pi,\pi\rbrack$ 上可积,$a_n,b_n$ 是它的 Fourier 系数,则:
    \begin{enumerate}
        \item 若 $f(-x)=f(x)$ 且 $f(\pi-x)=-f(x)$,则 $b_n=0$,且 $a_{2n}=0$;
        \item 若 $f(-x)=f(x)$ 且 $f(\pi-x)=f(x)$,则 $b_n=0$,且 $a_{2n-1}=0$;
        \item 若 $f(-x)=-f(x)$ 且 $f(\pi-x)=-f(x)$,则 $a_n=0$,且 $b_{2n-1}=0$;
        \item 若 $f(-x)=-f(x)$ 且 $f(\pi-x)=f(x)$,则 $a_n=0$,且 $b_{2n}=0$.
    \end{enumerate}
\end{proposition}

\begin{remark}
    $f(\pi-x)=-f(x)$ 表明图像关于点 $\left(\dfrac{\pi}{2},0\right)$ 中心对称,$f(\pi-x)=f(x)$ 表明图像关于直线 $x=\dfrac{\pi}{2}$ 轴对称.可用``偶心奇轴皆无偶,奇心偶轴皆无奇''记忆.
\end{remark}

%例题9.5.1
\begin{example}
    设
    \[
        \dfrac{a_0}{2}+\sum\limits_{k=1}^{\infty}(a_k\cos kx+b_k\sin kx)
    \]
    在 $\lbrack -\pi,\pi\rbrack$ 上一致收敛,试证它必是 $\lbrack -\pi,\pi\rbrack$ 上其和函数的 Fourier 级数.
\end{example}
\exampleblank

\begin{solution}
    一致收敛允许逐项积分.将和函数记为 $f$,分别乘以
    $1,\cos mx,\sin mx$ 后在 $\lbrack-\pi,\pi\rbrack$ 上积分,利用三角函数正交性,恰得
    \[
        a_m=\dfrac{1}{\pi}\int_{-\pi}^{\pi}f(x)\cos mx\,\mathrm{d}x,qquad
        b_m=\dfrac{1}{\pi}\int_{-\pi}^{\pi}f(x)\sin mx\,\mathrm{d}x.
    \]
    故原级数正是 $f$ 的 Fourier 级数.
\end{solution}

%例题9.5.2
\begin{example}[\markstar]
    三角多项式
    \[
        \dfrac{a_0}{2}+\sum\limits_{k=1}^{n}(a_k\cos kx+b_k\sin kx)
    \]
    的 Fourier 级数是它本身.
\end{example}
\exampleblank

\begin{solution}
    三角多项式与任意更高频率的正弦,余弦正交.按 Fourier 系数公式积分可知:次数不超过 $n$ 的系数保持原值,所有更高阶系数均为零.因此它的 Fourier 级数就是它本身.
\end{solution}

%例题9.5.3
\begin{example}
    设 $f(x)=\pi-x$, $x\in(0,\pi)$.
    \begin{enumerate}
        \item 将 $f(x)$ 展开为正弦级数;
        \item 判断该级数在 $(0,\pi)$ 内是否一致收敛.
    \end{enumerate}
\end{example}
\exampleblank

\begin{solution}
    正弦系数为
    \[
        b_n=\dfrac{2}{\pi}\int_0^\pi(\pi-x)\sin nx\,\mathrm{d}x=\dfrac{2}{n}.
    \]
    故
    \[
        \pi-x=2\sum\limits_{n=1}^{\infty}\dfrac{\sin nx}{n},qquad0<x<\pi.
    \]
    该级数不在 $(0,\pi)$ 一致收敛,因为每个部分和在 $x\to0^+$ 时趋于 $0$,而和函数趋于 $\pi$,故一致误差的下确界不小于 $\pi$.
\end{solution}

%例题9.5.4
\begin{example}
    试将
    \[
        f(x)=
        \begin{cases}
            1-x, & 0\leq x\leq2, \\
            x-3, & 2<x\leq4
        \end{cases}
    \]
    在 $\lbrack 0,4\rbrack$ 上展开为余弦级数,并证明
    \[
        \displaystyle\sum\limits_{n=0}^{\infty}\dfrac{1}{(2n+1)^2}=\dfrac{\pi^2}{8}.
    \]
\end{example}
\exampleblank

\begin{solution}
    余弦系数计算得 $a_0=0$,并且
    \[
        a_n=\dfrac{8}{\pi^2n^2}\left[(-1)^n-2\cos\dfrac{n\pi}{2}+1\right].
    \]
    故只有 $n=4k+2$ 时非零,
    \[
        f(x)=\dfrac{8}{\pi^2}\sum\limits_{k=0}^{\infty}
        \dfrac{\cos((4k+2)\pi x/4)}{(2k+1)^2}.
    \]
    令 $x=0$,由 $f(0)=1$ 得
    $\displaystyle\sum_{k=0}^{\infty}(2k+1)^{-2}=\pi^2/8$.
\end{solution}

%例题9.5.5
\begin{example}
    求
    \[
        f(x)=
        \begin{cases}
            x,   & 0\leq x<1,   \\
            2-x, & 1\leq x\leq2
        \end{cases}
    \]
    在 $\lbrack 0,2\rbrack$ 上的 Fourier 展开式.
\end{example}
\exampleblank

\begin{solution}
    按周期 $2$ 展开,计算得 $a_0=1$,$b_n=0$,且偶数阶余弦系数为零,奇数阶为
    $-4/(\pi^2n^2)$.因此
    \[
        f(x)=\dfrac{1}{2}-\dfrac{4}{\pi^2}
        \sum\limits_{k=0}^{\infty}\dfrac{\cos((2k+1)\pi x)}{(2k+1)^2}.
    \]
\end{solution}

%例题9.5.6
\begin{example}
    求
    \[
        f(x)=
        \begin{cases}
            0,      & -\pi\leq x\leq0, \\
            \sin x, & 0<x\leq\pi
        \end{cases}
    \]
    在 $\lbrack -\pi,\pi\rbrack$ 上的 Fourier 展开式.
\end{example}
\exampleblank

\begin{solution}
    系数为
    \[
        \dfrac{a_0}{2}=\dfrac{1}{\pi},\qquad b_1=\dfrac{1}{2},\quad b_n=0\ (n\ne1),
    \]
    以及
    \[
        a_{2k}=-\dfrac{2}{\pi(4k^2-1)},qquad a_{2k-1}=0.
    \]
    故
    \[
        f(x)=\dfrac{1}{\pi}+\dfrac{1}{2}\sin x-
        \dfrac{2}{\pi}\sum\limits_{k=1}^{\infty}\dfrac{\cos2kx}{4k^2-1},
    \]
    在跳跃点取左右极限的平均值.
\end{solution}

%例题9.5.7
\begin{example}
    设 $f(x)=x$, $x\in\lbrack 0,\dfrac{\pi}{2}\rbrack$,试将 $f(x)$ 展开为
    \[
        \displaystyle\sum\limits_{n=1}^{\infty}b_{2n-1}\sin(2n-1)x
    \]
    型的三角级数.
\end{example}
\exampleblank

\begin{solution}
    由正交性
    \[
        b_{2n-1}=\dfrac{4}{\pi}\int_0^{\pi/2}x\sin((2n-1)x)\,\mathrm{d}x
        =\dfrac{4(-1)^{n-1}}{\pi(2n-1)^2}.
    \]
    因此
    \[
        x=\dfrac{4}{\pi}\sum\limits_{n=1}^{\infty}
        \dfrac{(-1)^{n-1}}{(2n-1)^2}\sin((2n-1)x),qquad0\leq x\leq\dfrac{\pi}{2}.
    \]
\end{solution}

%例题9.5.8
\begin{example}
    已知 $f(x)=x$ 在 $(-\pi,\pi)$ 上的 Fourier 级数展开式为
    \[
        x=\sum\limits_{n=1}^{\infty}\dfrac{2(-1)^{n-1}}{n}\sin nx,
        \qquad |x|<\pi.
    \]
    试求函数 $\varphi(x)=x\sin x$ 在 $\lbrack -\pi,\pi\rbrack$ 上的 Fourier 展开式.
\end{example}
\exampleblank

\begin{solution}
    函数 $x\sin x$ 为偶函数,故只有余弦项.计算或把已知级数乘以 $\sin x$ 后用积化和差公式,得
    \[
        x\sin x=1-\dfrac{1}{2}\cos x-
        2\sum\limits_{n=2}^{\infty}\dfrac{(-1)^n}{n^2-1}\cos nx.
    \]
\end{solution}

%例题9.5.9
\begin{example}
    设 $f(x)$ 有 Fourier 系数
    \[
        f(x)\sim\dfrac{a_0}{2}+\sum\limits_{n=1}^{\infty}(a_n\cos nx+b_n\sin nx).
    \]
    试证逐项乘 $\sin x$ 可得 $f(x)\sin x$ 的 Fourier 级数.
\end{example}
\exampleblank

\begin{solution}
    用积化和差公式
    \[
        \cos nx\sin x=\dfrac{1}{2}[\sin(n+1)x-\sin(n-1)x],
    \]
    \[
        \sin nx\sin x=\dfrac{1}{2}[\cos(n-1)x-\cos(n+1)x].
    \]
    将有限部分和乘以 $\sin x$ 并按频率合并;再由 Fourier 系数积分公式取极限,便得到 $f(x)\sin x$ 的系数.等价地,令 $a_{-n}=a_n,b_{-n}=-b_n$,则新系数满足
    $A_n=(b_{n+1}-b_{n-1})/2$,$B_n=(a_{n-1}-a_{n+1})/2$.
\end{solution}

%例题9.5.10
\begin{example}[\markstar]
    \begin{enumerate}
        \item 设 $f(x)$ 以 $2\pi$ 为周期,在 $(0,2\pi)$ 内有界.若 $f(x)$ 单调递减,试证 $b_n\geq0$;若单调递增,则 $b_n\leq0$;
        \item 设 $f'(x)$ 在 $(0,2\pi)$ 内有界.若 $f'(x)$ 单调递增,试证 $a_n\geq0$;若单调递减,则 $a_n\leq0$;
        \item 设 $f(x)$ 在 $(0,2\pi)$ 内可积.若 $F(x)=\displaystyle\int_0^x\left(f(t)-\dfrac{a_0}{2}\right)\mathrm{d}t$ 单调递减,试证 $a_n\geq0$;若 $F$ 单调递增,则 $a_n\leq0$.
    \end{enumerate}
\end{example}
\exampleblank

\begin{solution}
    第 (1) 问把每个正半波与负半波配对:
    \[
        b_n=\dfrac{1}{\pi}\sum\int_0^{\pi/n}
        [f(t+2k\pi/n)-f(t+(2k+1)\pi/n)]\sin(nt)\,\mathrm{d}t.
    \]
    若 $f$ 递减,括号非负,故 $b_n\geq0$;递增时反号.

    第 (2) 问分部积分得 $a_n=-b_n(f')/n$,再用第 (1) 问.第 (3) 问因
    $F'=f-a_0/2$ 且 $F(0)=F(2\pi)=0$,分部积分得 $a_n=nb_n(F)$,再用第 (1) 问即可.
\end{solution}

\begin{proposition}[\markstar]
    若 $f$ 为周期 $2\pi$ 的可积和绝对可积函数,则其 Fourier 系数 $\{a_n\}$ 和 $\{b_n\}$ 必为无穷小量:
    \[
        \displaystyle\lim\limits_{n\to\infty}a_n=\lim\limits_{n\to\infty}b_n=0.
    \]
\end{proposition}

\begin{proposition}[\markstar]
    设以 $2\pi$ 为周期的函数 $f$ 在 $\lbrack -\pi,\pi\rbrack$ 上除了有限个点以外均有 $k+1$ 阶导数.如果 $f^{(k)}$ 在 $\lbrack -\pi,\pi\rbrack$ 上处处连续且 $f^{(k+1)}$ 可积和绝对可积,则
    \[
        a_n=o\left(\dfrac{1}{n^{k+1}}\right),
        \qquad
        b_n=o\left(\dfrac{1}{n^{k+1}}\right).
    \]
\end{proposition}

%例题9.5.11
\begin{example}
    设 $f$ 是以 $2\pi$ 为周期的函数且存在 $\alpha\in(0,1\rbrack$,使 $f$ 满足 $\alpha$ 阶 Lipschitz 条件
    \[
        |f(x)-f(y)|\leq L|x-y|^\alpha.
    \]
    证明:
    \[
        a_n=o\left(\dfrac{1}{n^\alpha}\right),
        \qquad
        b_n=o\left(\dfrac{1}{n^\alpha}\right).
    \]
\end{example}
\exampleblank

\begin{solution}
    题面在 $0<\alpha<1$ 时一般只能推出
    $a_n,b_n=O(n^{-\alpha})$,不能推出小 $o$.例如适当的 Weierstrass 型函数
    $\displaystyle\sum 2^{-k\alpha}\cos(2^kx)$ 属于 $\operatorname{Lip}\alpha$,但在 $n=2^k$ 处有
    $n^\alpha a_n$ 不趋于零.若把条件加强为小 Lipschitz 条件
    $\omega(f,h)=o(h^\alpha)$,则由平移恒等式
    \[
        (1-e^{inh})(a_n-ib_n)=\dfrac{1}{\pi}\int_{-\pi}^{\pi}[f(x)-f(x+h)]e^{-inx}\,\mathrm{d}x
    \]
    取 $h=\pi/n$ 即得题述小 $o$ 结论.$\alpha=1$ 时还可由几乎处处导数及 Riemann--Lebesgue 引理得到小 $o(1/n)$.
\end{solution}



\subsection{收敛定理}

\begin{theorem}[Fejer 定理 $*$]
    如果以 $2\pi$ 为周期的函数 $f$ 在 $\lbrack -\pi,\pi\rbrack$ 上可积和绝对可积,并且在点 $x_0$ 处有左,右极限 $f(x_0^+)$ 与 $f(x_0^-)$,则 $f$ 的 Fourier 级数在点 $x_0$ 处在 Cesaro 意义下收敛于
    \[
        \dfrac{1}{2}[f(x_0^+)+f(x_0^-)].
    \]
    特别地,当 $f$ 在 $x_0$ 连续时,Cesaro 和等于 $f(x_0)$.
\end{theorem}

\begin{theorem}[\markstar]
    设以 $2\pi$ 为周期的函数 $f$ 在 $\lbrack -\pi,\pi\rbrack$ 上可积和绝对可积,点 $x_0$ 是 $f$ 的连续点或第一类间断点.如果 $f$ 的 Fourier 级数在点 $x_0$ 处收敛,则它一定收敛于 $f(x_0)$ 或 $\dfrac{1}{2}[f(x_0^+)+f(x_0^-)]$.
\end{theorem}

\begin{theorem}[\markstar\ 唯一性定理]
    设 $f$ 为周期 $2\pi$ 的连续函数,且其 Fourier 系数均等于零,则 $f\equiv0$.
\end{theorem}

\begin{theorem}[$*$]
    如果 $f$ 是以 $2\pi$ 为周期的连续函数,则它的 Fourier 级数在 Cesaro 意义下一致收敛于 $f$.
\end{theorem}

\begin{remark}
    此定理也给出了 Weierstrass 三角多项式逼近定理中的三角多项式构造
    \[
        \sigma_n(x)=\dfrac{1}{n}\sum\limits_{k=1}^{n}S_k(x).
    \]
\end{remark}

%例题9.5.12
\begin{example}
    证明:若 $f(x)$ 是周期为 $2\pi$ 的连续函数,在 $\lbrack -\pi,\pi\rbrack$ 上分段光滑,则 $f(x)$ 的 Fourier 级数一致收敛于 $f(x)$.
\end{example}
\exampleblank

\begin{solution}
    周期连续且分段光滑的函数满足 Dini 条件.对 Dirichlet 积分写成
    \[
        S_n(x)-f(x)=\dfrac{1}{\pi}\int_0^\pi
        \dfrac{f(x+t)+f(x-t)-2f(x)}{2\sin(t/2)}
        \sin\left(n+\dfrac{1}{2}\right)t\,\mathrm{d}t.
    \]
    分段光滑性给出统一的 Lipschitz 控制,故小 $t$ 部分一致小;远离零点的部分由分部积分或 Riemann--Lebesgue 引理一致趋零.因此 $S_n\rightrightarrows f$.
\end{solution}

%例题9.5.13
\begin{example}
    构造两个以 $2\pi$ 为周期的连续函数,使得其 Fourier 级数在 $\lbrack 0,\pi\rbrack$ 上一致收敛于 $0$.
\end{example}
\exampleblank

\begin{solution}
    在 $\lbrack0,2\pi\rbrack$ 上取
    \[
        f_1(x)=\begin{cases}0, & 0\leq x\leq\pi, \\ \sin^2(x-\pi),&\pi<x<2\pi,
        \end{cases}
        \quad
        f_2(x)=\begin{cases}0, & 0\leq x\leq\pi, \\ \sin^4(x-\pi),&\pi<x<2\pi,
        \end{cases}
    \]
    并作 $2\pi$ 周期延拓.二者连续且分段光滑,故 Fourier 级数一致收敛于自身;在 $\lbrack0,\pi\rbrack$ 上两者都恒为零.
\end{solution}

%例题9.5.14
\begin{example}
    设 $f\in C\lbrack 0,1\rbrack$ 且 $f(0)=0$.若
    \[
        \displaystyle\int_0^1f(x)\sin(n\pi x)\,\mathrm{d}x=0,
    \]
    则 $f\equiv0$.
\end{example}
\exampleblank

\begin{solution}
    题面应理解为等式对所有正整数 $n$ 成立.正弦系 $\{\sin(n\pi x)\}$ 在 $L^2(0,1)$ 中完备,故所有正弦 Fourier 系数为零推出 $f=0$ 几乎处处.又 $f$ 连续,因此 $f\equiv0$.若题面只指某个固定的 $n$,结论显然不成立.
\end{solution}

%例题9.5.15
\begin{example}
    设 $f(x)$ 是以 $2\pi$ 为周期的函数,在 $\lbrack -\pi,\pi\rbrack$ 上可积.已知其 Fourier 级数部分和可表示为 Dirichlet 积分
    \[
        S_n(x)=\dfrac{1}{\pi}\int_{-\pi}^{\pi}f(x+t)
        \dfrac{\sin\left((n+\dfrac{1}{2})t\right)}{2\sin\dfrac{t}{2}}\,\mathrm{d}t,
    \]
    其中
    \[
        D_n(t)=\dfrac{\sin\left((n+\dfrac{1}{2})t\right)}{2\sin\dfrac{t}{2}}
        =\dfrac{1}{2}+\cos t+\cos2t+\cdots+\cos nt.
    \]
    又令 $\displaystyle\sigma_n(x)=\dfrac{1}{n}\sum\limits_{k=0}^{n-1}S_k(x)$.试证:
    \begin{enumerate}
        \item $\displaystyle\sum\limits_{k=0}^{n-1}D_k(x)=\dfrac{1}{2}\left(\dfrac{\sin\dfrac{nx}{2}}{\sin\dfrac{x}{2}}\right)^2$;
        \item $\displaystyle\dfrac{1}{2n\pi}\int_{-\pi}^{\pi}\left(\dfrac{\sin\dfrac{nx}{2}}{\sin\dfrac{x}{2}}\right)^2\mathrm{d}x=1$;
        \item 对任意 $\delta>0$,
              \[
                  \dfrac{1}{n\pi}\int_{\delta}^{\pi}\left(\dfrac{\sin\dfrac{nx}{2}}{\sin\dfrac{x}{2}}\right)^2\mathrm{d}x\to0,
                  \qquad n\to\infty;
              \]
        \item 若 $f(x)$ 是以 $2\pi$ 为周期的连续函数,则 $\sigma_n(x)\rightrightarrows f(x)$ 于 $\lbrack -\pi,\pi\rbrack$ 上.
    \end{enumerate}
\end{example}
\exampleblank

\begin{solution}
    第 (1) 问直接把 $D_k$ 写成余弦和并交换求和,得到 Fejér 核公式.第 (2) 问由正交性积分,常数项为 $1$.第 (3) 问在 $|x|\geq\delta$ 上有
    $|\sin(x/2)|\geq c_\delta>0$,故积分除以 $n$ 后为 $O(1/n)$.

    第 (4) 问写成
    \[
        \sigma_n(x)=\dfrac{1}{2n\pi}\int_{-\pi}^{\pi}f(x+t)
        \left(\dfrac{\sin(nt/2)}{\sin(t/2)}\right)^2\mathrm{d}t.
    \]
    核非负,总质量为 $1$,且质量集中到 $t=0$.结合 $f$ 的一致连续性,把积分分成 $|t|<\delta$ 与其补集,即得一致收敛.
\end{solution}

\begin{theorem}[\markstar\ Fourier 级数的逐项积分定理]
    设 $f$ 为周期 $2\pi$ 的函数,在 $\lbrack -\pi,\pi\rbrack$ 上可积和绝对可积,且
    \[
        f(x)\sim\dfrac{a_0}{2}+\sum\limits_{n=1}^{\infty}(a_n\cos nx+b_n\sin nx).
    \]
    则 $\displaystyle\sum\limits_{n=1}^{\infty}\dfrac{b_n}{n}$ 收敛,并且
    \[
        \displaystyle\int_a^bf(t)\,\mathrm{d}t
        =\int_a^b\dfrac{a_0}{2}\,\mathrm{d}t
        +\sum\limits_{n=1}^{\infty}\int_a^b(a_n\cos nt+b_n\sin nt)\,\mathrm{d}t.
    \]
    对于 $a=0,b=x$ 可得
    \[
        \displaystyle\int_0^x\left(f(t)-\dfrac{a_0}{2}\right)\mathrm{d}t
        =\sum\limits_{n=1}^{\infty}\dfrac{b_n}{n}
        +\sum\limits_{n=1}^{\infty}\left(-\dfrac{b_n}{n}\cos nx+\dfrac{a_n}{n}\sin nx\right).
    \]
\end{theorem}

\begin{remark}
    由此可知一个三角级数绝对可积的必要条件为 $\displaystyle\sum\limits_{n=1}^{\infty}\dfrac{b_n}{n}$ 收敛.该结论区别于一般函数项级数情形:不论 $f$ 的 Fourier 级数收敛情况如何,都可以逐项积分.
\end{remark}

\begin{theorem}[\markstar\ Fourier 级数的逐项求导定理]
    设 $f$ 是以 $2\pi$ 为周期的连续函数,在 $\lbrack -\pi,\pi\rbrack$ 上除有限个点外处处可导,又设
    \[
        f(x)\sim\dfrac{a_0}{2}+\sum\limits_{n=1}^{\infty}(a_n\cos nx+b_n\sin nx),
    \]
    且 $f'(x)$ 在 $\lbrack -\pi,\pi\rbrack$ 上分段光滑,则
    \[
        \dfrac{f'(x^+)+f'(x^-)}{2}
        =\sum\limits_{n=1}^{\infty}(a_n\cos nx+b_n\sin nx)'
        =\sum\limits_{n=1}^{\infty}(nb_n\cos nx-na_n\sin nx).
    \]
\end{theorem}

%例题9.5.16
\begin{example}
    求
    \[
        f(x)=\ln(1+2r\cos x+r^2),
        \qquad |r|<1,
    \]
    的 Fourier 级数展开式.
\end{example}
\exampleblank

\begin{solution}
    因
    \[
        1+2r\cos x+r^2=|1+re^{ix}|^2,
    \]
    且 $|r|<1$,展开对数得
    \[
        \ln(1+2r\cos x+r^2)
        =2\Ree\sum\limits_{n=1}^{\infty}\dfrac{(-1)^{n+1}r^ne^{inx}}n
        =2\sum\limits_{n=1}^{\infty}\dfrac{(-1)^{n+1}r^n}{n}\cos nx.
    \]
\end{solution}

%例题9.5.17
\begin{example}
    已知
    \[
        x=2\sum\limits_{n=1}^{\infty}\dfrac{(-1)^{n+1}\sin nx}{n},
        \qquad -\pi<x<\pi.
    \]
    试求函数 $x^2,x^3$ 的 Fourier 级数展开式.
\end{example}
\exampleblank

\begin{solution}
    对已知级数逐项积分并定常数,得到
    \[
        x^2=\dfrac{\pi^2}{3}+4\sum\limits_{n=1}^{\infty}
        \dfrac{(-1)^n}{n^2}\cos nx.
    \]
    再计算奇函数 $x^3$ 的正弦系数,得
    \[
        x^3=2\sum\limits_{n=1}^{\infty}(-1)^n
        \left(-\dfrac{\pi^2}{n}+\dfrac{6}{n^3}\right)\sin nx,qquad |x|<\pi.
    \]
\end{solution}

%例题9.5.18
\begin{example}
    \begin{enumerate}
        \item 证明 $\displaystyle\sum\limits_{n=2}^{\infty}\dfrac{1}{\ln n}\sin nx$ 不可能为某函数的 Fourier 级数;
        \item 证明 $\displaystyle\sum\limits_{n=9}^{\infty}\dfrac{\sin nx}{\ln\ln n}$ 不可能为某函数的 Fourier 级数.
    \end{enumerate}
\end{example}
\exampleblank

\begin{solution}
    使用单调正弦级数的 Abel 型渐近公式:若 $c_n$ 正,单调趋零且缓慢变化,则
    \[
        \displaystyle\sum\limits_{n=1}^{\infty}c_n\sin(nx)\asymp\dfrac{c_{\lfloor1/x\rfloor}}x,qquad x\to0^+.
    \]
    对 $c_n=1/\ln n$,右端同阶于 $1/[x\ln(1/x)]$;对
    $c_n=1/\ln\ln n$,同阶于 $1/[x\ln\ln(1/x)]$.二者在零点附近都不可积.若它们是某个可积函数的 Fourier 级数,则 Abel 和应给出一个局部可积函数,矛盾.因此两级数都不可能是可积函数的 Fourier 级数.
\end{solution}



\subsection{Parseval 等式}

%例题9.5.19
\begin{example}
    求函数
    \[
        f(x)=\left(x-\dfrac{\pi}{2}\right)^2,
        \qquad x\in\lbrack 0,2\pi\rbrack,
    \]
    的 Fourier 级数,并求级数 $\displaystyle\sum\limits_{n=1}^{\infty}\dfrac{1}{n^4}$ 的和.
\end{example}
\exampleblank

\begin{solution}
    直接计算系数得
    \[
        f(x)\sim\dfrac{7\pi^2}{12}+
        \sum\limits_{n=1}^{\infty}\left(\dfrac{4}{n^2}\cos nx-\dfrac{2\pi}{n}\sin nx\right).
    \]
    由 Parseval 等式,
    \[
        \dfrac{1}{\pi}\int_0^{2\pi}f^2(x)\,\mathrm{d}x
        =\dfrac{1}{2}\left(\dfrac{7\pi^2}{6}\right)^2
        +\sum\limits_{n=1}^{\infty}\left(\dfrac{16}{n^4}+\dfrac{4\pi^2}{n^2}\right).
    \]
    左端为 $61\pi^4/40$,再用 $\displaystyle\sum n^{-2}=\pi^2/6$,解得
    $\displaystyle\sum n^{-4}=\pi^4/90$.
\end{solution}

%例题9.5.20
\begin{example}
    \begin{enumerate}
        \item 设 $f(x)$ 在 $\lbrack 0,\pi\rbrack$ 上有连续导数,且 $f'(x)$ 在 $\lbrack 0,\pi\rbrack$ 上分段光滑,$\displaystyle\int_0^\pi f(x)\,\mathrm{d}x=0$.证明相应 Wirtinger 型不等式;
        \item 设 $f(x)$ 在 $\mathbb{R}$ 上有连续导数,周期为 $l>0$,且 $\displaystyle\int_0^lf(x)\,\mathrm{d}x=0$.证明
              \[
                  \displaystyle\int_0^l|f'(x)|^2\,\mathrm{d}x
                  \geq\dfrac{4\pi^2}{l^2}\int_0^l|f(x)|^2\,\mathrm{d}x.
              \]
              等号成立当且仅当
              \[
                  f(x)=a_1\cos\dfrac{2\pi x}{l}+b_1\sin\dfrac{2\pi x}{l}.
              \]
    \end{enumerate}
\end{example}
\exampleblank

\begin{solution}
    第 (1) 问把 $f$ 作偶延拓并展开为余弦级数.零平均条件使常数项为零.Parseval 等式给出
    \[
        \displaystyle\int_0^\pi|f|^2\,\mathrm{d}x=\dfrac{\pi}{2}\sum\limits_{n\geq1}a_n^2,qquad
        \int_0^\pi|f'|^2\,\mathrm{d}x=\dfrac{\pi}{2}\sum\limits_{n\geq1}n^2a_n^2,
    \]
    故后一积分不小于前者;等号当且仅当只有 $n=1$ 项.

    第 (2) 问令 $t=2\pi x/l$,化为 $2\pi$ 周期情形.零平均消去常数项,Parseval 等式给出
    \[
        \displaystyle\int_0^l|f'|^2\,\mathrm{d}x\geq\dfrac{4\pi^2}{l^2}
        \int_0^l|f|^2\,\mathrm{d}x.
    \]
    等号当且仅当 Fourier 系数只含第一谐波,即题中所列形式.
\end{solution}
